\documentclass[11pt,a4paper]{article}
\usepackage[utf8]{inputenc}
\usepackage[T1]{fontenc}
\usepackage[ngerman]{babel}
\usepackage{lmodern}
\usepackage{microtype}
\usepackage[a4paper,margin=2.5cm]{geometry}
\usepackage{amsmath,amssymb}
\usepackage{booktabs,array,longtable}
\usepackage[table]{xcolor}
\usepackage{tikz}
\usepackage{pdflscape}
\usepackage{hyperref}
\usepackage{csquotes}
\hypersetup{colorlinks=true,linkcolor=blue,citecolor=blue,urlcolor=blue}
\usetikzlibrary{patterns}

\title{Auf dem Weg zu einem Antimaterie-Periodensystem:\\
Ein epistemisches Preparedness Framework fuer Anti-Atome, Anti-Ionen, Anti-Molekuele und kosmische Antimaterie-Domaenen}
\author{Emanuel und Dr. Martin Dres}
\date{Stand: 25. April 2026}

\newcommand{\anti}[1]{\ensuremath{\overline{\mathrm{#1}}}}
\newcommand{\ap}{\ensuremath{\bar p}}
\newcommand{\ep}{\ensuremath{e^+}}
\newcommand{\Zanti}{\ensuremath{Z_{\bar p}}}

\begin{document}
\maketitle

\begin{abstract}
Das klassische Periodensystem ordnet chemische Elemente nach Protonenzahl, Elektronenkonfiguration und daraus folgender Periodizitaet. Fuer Antimaterie gibt es bisher keine IUPAC-kodifizierte Reihe eigenstaendiger Antielement-Namen und kein etabliertes Periodensystem, das den empirischen Status einzelner Anti-Atome explizit sichtbar macht. Diese Arbeit schlaegt deshalb kein alternatives Periodensystem mit anderer Chemie vor, sondern ein \emph{epistemisches Preparedness Framework}: Die durch CPT-Symmetrie erwartete strukturelle Isomorphie wird von der experimentellen Evidenz getrennt. Antiwasserstoff ist der einzige neutrale Anti-Atomfall, der bisher systematisch erzeugt, gefangen und spektroskopisch sowie gravitativ untersucht wurde. Antihelium ist dagegen experimentell bisher als Antikern, nicht als neutrales Antihelium-Atom etabliert. Fuer hoehere Anti-Elemente ist die Orbital- und Periodizitaetsstruktur theoretisch erwartbar, waehrend neutrale Anti-Atome, Anti-Ionen, Anti-Molekuele, chemische Reaktionsdaten und moegliche kosmische Antimaterie-Domaenen offen bleiben. Der Beitrag des vorgeschlagenen Antiperiodensystems liegt daher in einer expliziten Status-Codierung: gesicherte Anti-Atome, beobachtete Antikerne, CPT-erwartete, aber nicht realisierte Anti-Atome, offene Chemie und spekulative astrophysikalische Domaenen.
\end{abstract}

\tableofcontents
\newpage

\section{Ziel und Abgrenzung}

Der Begriff \enquote{Antiperiodensystem} kann leicht missverstanden werden. Er darf nicht bedeuten, dass Antimaterie nach heutigem Standardmodell eine andere chemische Periodizitaet besitzen sollte. Unter den Voraussetzungen lokaler, Lorentz-invarianter und unitaerer Quantenfeldtheorie folgt aus dem CPT-Theorem, dass ein isoliertes Anti-Atom dieselben inneren Energieniveaus wie das entsprechende Atom besitzen sollte, sofern alle Ladungen, Paritaeten und Zeitrichtungen korrekt transformiert werden. Das vorgeschlagene System ist daher nicht als Konkurrenz zum Periodensystem der Materie gemeint, sondern als Forschungslandkarte.

Die Kernfrage dieser Arbeit lautet:

\begin{quote}
Wie kann man ein Periodensystem der Antimaterie so darstellen, dass gleichzeitig sichtbar wird, was physikalisch stark erwartet, was experimentell bestaetigt und was fuer komplexere Anti-Systeme noch offen ist?
\end{quote}

Daraus folgt eine klare Abgrenzung:

\begin{itemize}
  \item \textbf{Nicht neu} ist die Idee, dass Antiatome ladungskonjugierte Gegenstuecke normaler Atome sind.
  \item \textbf{Nicht behauptet} wird eine alternative Antimaterie-Chemie im Sinne anderer Valenzen oder anderer Gruppenstruktur.
  \item \textbf{Vorgeschlagen} wird eine explizite, statuscodierte Tabelle, die Antinukleus-, Anti-Atom-, Spektroskopie-, Chemie- und Astrophysikstatus unterscheidet.
\end{itemize}

Diese Trennung ist wichtig, weil die bisherige Forschung terminologisch und experimentell stark um Antiwasserstoff zentriert ist. CERN beschreibt Antiwasserstoff als neutrales Anti-Atom aus Antiproton und Positron; ALPHA, ATRAP, AEGIS, GBAR, ASACUSA und verwandte Experimente nutzen ihn zur Pruefung von Materie-Antimaterie-Symmetrien \cite{CERNstoring,CERNantimatter}.

\section{Nomenklatur: Was ist vergeben, was ist nur Konvention?}

\subsection{Befund}

Nach der hier herangezogenen IUPAC- und Fachliteratur gibt es keine eigenstaendige, von IUPAC ratifizierte Namensreihe fuer \enquote{Antielemente}. IUPAC definiert klassische Element- und Isotopennamen im Rahmen normaler Atome und Nuclide; die Ordnungszahl ist die Protonenzahl \cite{IUPACatomicnumber}, und die offiziellen Elementnamen und Symbole werden fuer Elemente festgelegt, nicht fuer eine separate Antielement-Reihe \cite{IUPAC113118}. Fuer Antimaterie wird in der Physik stattdessen eine kombinierte Konvention verwendet: der Praefix \enquote{anti-} im Namen und ein Ueberstrich im Symbol, etwa \anti{H} fuer Antiwasserstoff.

IUPAC selbst enthaelt relevante Einzelbegriffe wie \enquote{Positron} als positiv geladenes Elektron \cite{IUPACpositron} und \enquote{Annihilation} als Wechselwirkung eines Teilchens mit seinem Antiteilchen \cite{IUPACannihilation}. Das ist aber keine Antielement-Nomenklatur. Daraus folgt: Ein Antiperiodensystem sollte keine neuen offiziellen Namen erfinden, sondern transparent als Arbeitsnotation arbeiten.

\subsection{Empfohlene Terminologie fuer das Paper}

\begin{longtable}{p{0.22\textwidth}p{0.24\textwidth}p{0.18\textwidth}p{0.27\textwidth}}
\caption{Terminologische Trennung fuer ein Antimaterie-Periodensystem.}\label{tab:nomenclature}\\
\toprule
\textbf{Gegenstand} & \textbf{Empfohlene Bezeichnung} & \textbf{Symbol/Notation} & \textbf{Status und Kommentar}\\
\midrule
\endfirsthead
\toprule
\textbf{Gegenstand} & \textbf{Empfohlene Bezeichnung} & \textbf{Symbol/Notation} & \textbf{Status und Kommentar}\\
\midrule
\endhead
Antiproton & Antiproton & \ap & Etablierte Teilchenphysik-Notation; kein chemisches Element.\\
Positron & Positron / Antielektron & \ep & IUPAC definiert das Positron als positiv geladenes Elektron \cite{IUPACpositron}.\\
Neutrales Antiwasserstoffatom & Antiwasserstoff & \anti{H} & Etablierter Begriff in CERN/ALPHA-Literatur; besteht aus \ap{} und \ep{} \cite{CERNstoring}.\\
Antiwasserstoff-Ion & Antiwasserstoff-Ion & \anti{H}$^{+}$ & GBAR erzeugt bzw. nutzt positiv geladene Antiwasserstoff-Ionen als Zwischenschritt fuer Gravitationsexperimente \cite{CERNGBAR}.\\
Antihelium-Kern & Antihelium-3- oder Antihelium-4-Kern & $\overline{{}^{3}\mathrm{He}}$, $\overline{{}^{4}\mathrm{He}}$ & Experimentell als Antikern bzw. Anti-$\alpha$ beobachtet; nicht gleichbedeutend mit neutralem Antihelium-Atom \cite{STAR2011}.\\
Antiprotonisches Helium & antiprotonisches Helium & $\bar p\mathrm{He}^{+}$ u. a. & Exotisches Atom aus normalem Heliumkern, Elektron und Antiproton; nicht Antihelium \cite{Hori2011}.\\
Allgemeines Anti-Element & Anti- + Elementname & \anti{E} bzw. \anti{X} & Sinnvolle Arbeitsnotation; keine IUPAC-Neubenennung. Fuer mehrbuchstabige Symbole sollte der Ueberstrich ueber dem ganzen Symbol stehen: \anti{He}, nicht nur $\bar H e$.\\
\bottomrule
\end{longtable}

Fuer dieses Manuskript wird daher folgende Regel verwendet:

\begin{equation}
E_Z \longmapsto \overline{E}_Z,
\end{equation}

wobei $Z$ weiterhin der positive Index der Position im klassischen Periodensystem bleibt. Fuer Antimaterie bezeichnet $Z$ operational die Anzahl der Antiprotonen im Antikern:

\begin{equation}
\Zanti = N_{\bar p} = Z.
\end{equation}

Die Kernladung des Antikerns ist zwar $-Ze$, aber das Periodensystem sollte nicht mit negativen oder ueberstrichenen Ordnungszahlen arbeiten. Sonst wuerde die Ordnungsrelation kuenstlich gebrochen. Die Position des Anti-Elements wird durch die Materie-Position $Z$ indexiert; nur die Konstituenten sind ladungskonjugiert.

\section{Stand des Wissens}

\subsection{Antiwasserstoff}

Antiwasserstoff ist der zentrale experimentelle Anker. CERN beschreibt Antiwasserstoff als neutralen Anti-Atomzustand aus einem negativ geladenen Antiproton und einem positiv geladenen Positron \cite{CERNstoring}. ALPHA hat Antiwasserstoff gefangen, laser-spektroskopisch untersucht und die $1S$-$2S$-Transition mit hoher Praezision charakterisiert \cite{ALPHA2018}. Die ALPHA-g-Messung von 2023 zeigte, dass freigesetzte Antiwasserstoffatome in einer Weise fallen, die mit gravitativer Anziehung zur Erde vereinbar ist; repulsive \enquote{Antigravitation} fuer diesen Fall wurde ausgeschlossen \cite{Anderson2023}. 2025 berichtete ALPHA zudem eine deutlich verbesserte simultane Speicherung von mehr als 15.000 Antiwasserstoffatomen; dieselbe Arbeit bezeichnet Antiwasserstoff als das einzige reine antiatomare System, das bisher untersucht wurde \cite{Akbari2025}.

\subsection{Anti-Ionen, Anti-Molekuele und Antimaterie-Chemie}

Die neue Literatur zeigt, dass der Uebergang von isoliertem Antiwasserstoff zu Antimaterie-Chemie bereits vorbereitet wird. Zammit et al. veroeffentlichten 2025 eine systematische Perspektive zur \enquote{Antihydrogen chemistry}. Darin werden Reaktionswege zur Bildung von Antiwasserstoff sowie anionischen, kationischen und molekularen Antimaterie-Spezies durch kollisions- und strahlungsgetriebene Prozesse diskutiert \cite{Zammit2025}. Das stuetzt die Grundidee dieses Papers: Ein Antiperiodensystem ist heute weniger eine Tabelle gemessener chemischer Daten als eine strukturierte Forschungs- und Vorhersagekarte.

\subsection{Antihelium und schwere Antimaterie-Kernsysteme}

Antihelium ist nomenklatorisch besonders wichtig. Die STAR Collaboration beobachtete 2011 den Antihelium-4-Kern, also das antimaterielle Gegenstueck des $\alpha$-Teilchens, bestehend aus zwei Antiprotonen und zwei Antineutronen \cite{STAR2011}. Das ist ein Antikernnachweis, aber kein Nachweis eines neutralen Antihelium-Atoms mit Positronenhuelle. 2024 berichtete STAR zudem ein Antimaterie-Hypernukleus-System \cite{STAR2024}. Solche Ergebnisse zeigen, dass Kern-Antimaterie komplexer als Antiwasserstoff erzeugt werden kann; sie begruenden aber noch keine Anti-Atom- oder Antimolekuelchemie fuer hoehere Elemente.

\subsection{Kosmische Antimaterie-Domaenen}

Makroskopische Antimaterie-Domaenen, Anti-Sterne oder Antimaterie-Cluster sind fuer dieses Framework ein eigener Status-Layer, nicht Teil der gesicherten Chemie. CERN beschreibt, dass Antimaterie auf der Erde nur in sehr kleinen Mengen hergestellt oder als Einzelteilchen in natuerlichen Prozessen gefunden wird; die grosse Materie-Antimaterie-Asymmetrie bleibt offen \cite{CERNantimatter}. Antikerne in kosmischer Strahlung waeren astrophysikalisch hoch relevant. Das Antiperiodensystem sollte deshalb eine Spalte oder Markierung fuer \enquote{astro-physikalisch gesucht / unbestaetigt} enthalten, ohne daraus eine chemische Eigenschaft abzuleiten.

\section{Physikalisches Strukturmodell}

\subsection{CPT-isomorphe Abbildung}

Ein neutrales Atom mit Kernladung $+Ze$ und $Z$ Elektronen wird formal auf ein neutrales Anti-Atom mit Antikernladung $-Ze$ und $Z$ Positronen abgebildet:

\begin{equation}
\mathcal{A}_Z = \left(N_Z, e^-_1,\ldots,e^-_Z\right)
\quad \xrightarrow{CPT} \quad
\overline{\mathcal{A}}_Z = \left(\overline{N}_Z, e^+_1,\ldots,e^+_Z\right).
\end{equation}

Im nichtrelativistischen Coulomb-Anteil des Hamiltonoperators bleiben die relevanten Ladungsprodukte invariant. Fuer die Kern-Huellen-Anziehung gilt:

\begin{equation}
(+Ze)(-e)=(-Ze)(+e)=-Ze^2,
\end{equation}

und fuer die Lepton-Lepton-Abstossung:

\begin{equation}
(-e)(-e)=(+e)(+e)=+e^2.
\end{equation}

Damit ist die interne Coulomb-Struktur eines isolierten Anti-Atoms formal isomorph zur Materie-Struktur. Relativistische, QED- und finite-Kerngroessen-Effekte werden nicht ignoriert; sie werden ebenfalls ladungskonjugiert. Unter CPT sollten die Eigenwerte uebereinstimmen. Eine gemessene Abweichung waere daher keine normale chemische Variante, sondern ein Hinweis auf neue Physik oder auf nicht beruecksichtigte Umgebungsbedingungen.

\subsection{Was CP-Verletzung und Gravitation bedeuten - und was nicht}

CP-Verletzung ist in schwachen Wechselwirkungen bekannt, impliziert aber nicht automatisch andere chemische Valenzen von Antiatomen. Gewoehnliche chemische Bindung wird primaer durch elektromagnetische Wechselwirkung und die Quantenzustandsstruktur der Leptonenhuelle bestimmt. Fuer ein Antiperiodensystem bedeutet das:

\begin{itemize}
  \item \textbf{Chemische Periodizitaet:} Nach Standardphysik identisch erwartbar.
  \item \textbf{CPT-Tests:} Praezisionsspektroskopie und Gravitationsexperimente koennen Abweichungen suchen.
  \item \textbf{Komplexe Systeme:} Anti-Molekuele, Anti-Ionen und Antimaterie-Cluster koennen experimentell neue Herausforderungen zeigen, ohne dass die Tabellenstruktur selbst anders sein muss.
  \item \textbf{Astrophysik:} Kosmische Antimaterie-Domaenen waeren ein kosmologisches Problem der Materie-Antimaterie-Asymmetrie, nicht zunaechst ein anderes Periodengesetz.
\end{itemize}

\section{Vorschlag: epistemisches Antiperiodensystem}

\subsection{Warum ein Skalar nicht reicht}

Der urspruengliche Confidence Layer $C(Z)$ ist als einfache Visualisierung nuetzlich, aber wissenschaftlich zu grob. Ein einzelner Wert vermischt unterschiedliche Fragen:

\begin{enumerate}
  \item Ist die Bezeichnung etabliert?
  \item Ist ein Antikern beobachtet?
  \item Ist ein neutrales Anti-Atom erzeugt?
  \item Gibt es Spektroskopie?
  \item Gibt es Chemie oder Molekuele?
  \item Gibt es astrophysikalische Evidenz?
\end{enumerate}

Deshalb wird hier ein Vektorstatus vorgeschlagen:

\begin{equation}
\mathbf{S}(\overline{E}_Z)=
\left(S_{\mathrm{name}},S_{\mathrm{nuc}},S_{\mathrm{atom}},S_{\mathrm{spec}},S_{\mathrm{chem}},S_{\mathrm{astro}}\right).
\end{equation}

Jede Komponente kann mit vier Stufen codiert werden:

\begin{center}
\begin{tabular}{clp{0.55\textwidth}}
\toprule
\textbf{Stufe} & \textbf{Name} & \textbf{Bedeutung}\\
\midrule
0 & nicht definiert / nicht belegt & Keine belastbare Terminologie oder Evidenz.\\
1 & theoretisch erwartet & Aus CPT, QED und bekannter Materie-Chemie abgeleitet, aber nicht experimentell realisiert.\\
2 & beobachtet & Antikern, Anti-Atom oder Spezies experimentell erzeugt oder nachgewiesen.\\
3 & praezise charakterisiert & Spektroskopie, Gravitation, Reaktionsrate oder andere quantitative Daten liegen vor.\\
\bottomrule
\end{tabular}
\end{center}

\subsection{Minimaler Status nach heutigem Stand}

\begin{center}
\begin{tabular}{lcccccc}
\toprule
\textbf{Gegenstand} & $S_{\mathrm{name}}$ & $S_{\mathrm{nuc}}$ & $S_{\mathrm{atom}}$ & $S_{\mathrm{spec}}$ & $S_{\mathrm{chem}}$ & $S_{\mathrm{astro}}$\\
\midrule
\anti{H} & 3 & 3 & 3 & 3 & 1--2 & 0\\
\anti{He} & 2 & 2 & 0 & 0 & 0--1 & 0\\
\anti{Li} bis \anti{Ne} & 1 & 0--1 & 0 & 0 & 0--1 & 0\\
\anti{Na} bis \anti{Og} & 1 & 0 & 0 & 0 & 0--1 & 0\\
Kosmische Antimaterie-Domaenen & 1 & 0--1 & 0 & 0 & 0 & 0\\
\bottomrule
\end{tabular}
\end{center}

Die Tabelle ist bewusst vorsichtig. Sie unterscheidet zwischen struktureller Erwartung und experimenteller Realisierung. Antihelium hat Kernstatus, aber keinen neutralen Anti-Atomstatus. Antiwasserstoff hat den klar hoechsten Status. Fuer Antimaterie-Chemie ausserhalb wasserstoffnaher Systeme ist der Status ueberwiegend vorbereitet, nicht gemessen.

\subsection{Visualisierungsprinzip}

Die beste Darstellung ist eine Split-Cell-Darstellung:

\begin{itemize}
  \item Die obere/rechte Zellhaelfte zeigt die klassische Orbitalblock-Struktur der Materie: s-, p-, d- und f-Block.
  \item Die untere/linke Zellhaelfte zeigt den Status des Antimaterie-Gegenstuecks.
  \item Der Ueberstrich liegt nur auf dem Elementsymbol, nicht auf der Ordnungszahl.
  \item Die Farbcodierung sagt nicht: \enquote{die Theorie ist unsicher}, sondern: \enquote{der experimentelle Realisierungs- und Charakterisierungsstand ist unterschiedlich}.
\end{itemize}

\begin{figure}[H]
\centering
\begin{tikzpicture}[x=0.68cm, y=0.68cm, every node/.style={inner sep=0pt}]
\path[use as bounding box] (0,0.1) rectangle (18,-12.2);
\fill[blue!18, rounded corners=0.7pt] (0,0) -- (1,0) -- (1,-1) -- cycle;
\fill[green!60, rounded corners=0.7pt] (0,0) -- (0,-1) -- (1,-1) -- cycle;
\draw[black, rounded corners=0.7pt] (0,0) rectangle (1,-1);
\draw[black, dashed, very thin] (0,0) -- (1,-1);
\node at (0.7,-0.25) {\tiny \textbf{H}};
\node at (0.28,-0.73) {\tiny \textbf{$\overline{\mathrm{H}}$}};
\node at (0.18,-0.18) {\tiny 1};
\node[fill=white, draw=black, circle, inner sep=0.4pt] at (0.8,-0.8) {\tiny \textbf{=}};
\fill[green!18, rounded corners=0.7pt] (17,0) -- (18,0) -- (18,-1) -- cycle;
\fill[orange!60, rounded corners=0.7pt] (17,0) -- (17,-1) -- (18,-1) -- cycle;
\draw[black, rounded corners=0.7pt] (17,0) rectangle (18,-1);
\draw[black, dashed, very thin] (17,0) -- (18,-1);
\node at (17.7,-0.25) {\tiny \textbf{He}};
\node at (17.28,-0.73) {\tiny \textbf{$\overline{\mathrm{He}}$}};
\node at (17.18,-0.18) {\tiny 2};
\node[fill=white, draw=black, circle, inner sep=0.4pt] at (17.8,-0.8) {\tiny \textbf{K}};
\fill[blue!18, rounded corners=0.7pt] (0,-1) -- (1,-1) -- (1,-2) -- cycle;
\fill[gray!8, rounded corners=0.7pt] (0,-1) -- (0,-2) -- (1,-2) -- cycle;
\draw[black, rounded corners=0.7pt] (0,-1) rectangle (1,-2);
\draw[black, dashed, very thin] (0,-1) -- (1,-2);
\node at (0.7,-1.25) {\tiny \textbf{Li}};
\node at (0.28,-1.73) {\tiny \textbf{$\overline{\mathrm{Li}}$}};
\node at (0.18,-1.18) {\tiny 3};
\node[fill=white, draw=black, circle, inner sep=0.4pt] at (0.8,-1.8) {\tiny \textbf{?}};
\fill[blue!18, rounded corners=0.7pt] (1,-1) -- (2,-1) -- (2,-2) -- cycle;
\fill[gray!8, rounded corners=0.7pt] (1,-1) -- (1,-2) -- (2,-2) -- cycle;
\draw[black, rounded corners=0.7pt] (1,-1) rectangle (2,-2);
\draw[black, dashed, very thin] (1,-1) -- (2,-2);
\node at (1.7,-1.25) {\tiny \textbf{Be}};
\node at (1.28,-1.73) {\tiny \textbf{$\overline{\mathrm{Be}}$}};
\node at (1.18,-1.18) {\tiny 4};
\node[fill=white, draw=black, circle, inner sep=0.4pt] at (1.8,-1.8) {\tiny \textbf{?}};
\fill[green!18, rounded corners=0.7pt] (12,-1) -- (13,-1) -- (13,-2) -- cycle;
\fill[gray!8, rounded corners=0.7pt] (12,-1) -- (12,-2) -- (13,-2) -- cycle;
\draw[black, rounded corners=0.7pt] (12,-1) rectangle (13,-2);
\draw[black, dashed, very thin] (12,-1) -- (13,-2);
\node at (12.7,-1.25) {\tiny \textbf{B}};
\node at (12.28,-1.73) {\tiny \textbf{$\overline{\mathrm{B}}$}};
\node at (12.18,-1.18) {\tiny 5};
\node[fill=white, draw=black, circle, inner sep=0.4pt] at (12.8,-1.8) {\tiny \textbf{?}};
\fill[green!18, rounded corners=0.7pt] (13,-1) -- (14,-1) -- (14,-2) -- cycle;
\fill[gray!8, rounded corners=0.7pt] (13,-1) -- (13,-2) -- (14,-2) -- cycle;
\draw[black, rounded corners=0.7pt] (13,-1) rectangle (14,-2);
\draw[black, dashed, very thin] (13,-1) -- (14,-2);
\node at (13.7,-1.25) {\tiny \textbf{C}};
\node at (13.28,-1.73) {\tiny \textbf{$\overline{\mathrm{C}}$}};
\node at (13.18,-1.18) {\tiny 6};
\node[fill=white, draw=black, circle, inner sep=0.4pt] at (13.8,-1.8) {\tiny \textbf{?}};
\fill[green!18, rounded corners=0.7pt] (14,-1) -- (15,-1) -- (15,-2) -- cycle;
\fill[gray!8, rounded corners=0.7pt] (14,-1) -- (14,-2) -- (15,-2) -- cycle;
\draw[black, rounded corners=0.7pt] (14,-1) rectangle (15,-2);
\draw[black, dashed, very thin] (14,-1) -- (15,-2);
\node at (14.7,-1.25) {\tiny \textbf{N}};
\node at (14.28,-1.73) {\tiny \textbf{$\overline{\mathrm{N}}$}};
\node at (14.18,-1.18) {\tiny 7};
\node[fill=white, draw=black, circle, inner sep=0.4pt] at (14.8,-1.8) {\tiny \textbf{?}};
\fill[green!18, rounded corners=0.7pt] (15,-1) -- (16,-1) -- (16,-2) -- cycle;
\fill[gray!8, rounded corners=0.7pt] (15,-1) -- (15,-2) -- (16,-2) -- cycle;
\draw[black, rounded corners=0.7pt] (15,-1) rectangle (16,-2);
\draw[black, dashed, very thin] (15,-1) -- (16,-2);
\node at (15.7,-1.25) {\tiny \textbf{O}};
\node at (15.28,-1.73) {\tiny \textbf{$\overline{\mathrm{O}}$}};
\node at (15.18,-1.18) {\tiny 8};
\node[fill=white, draw=black, circle, inner sep=0.4pt] at (15.8,-1.8) {\tiny \textbf{?}};
\fill[green!18, rounded corners=0.7pt] (16,-1) -- (17,-1) -- (17,-2) -- cycle;
\fill[gray!8, rounded corners=0.7pt] (16,-1) -- (16,-2) -- (17,-2) -- cycle;
\draw[black, rounded corners=0.7pt] (16,-1) rectangle (17,-2);
\draw[black, dashed, very thin] (16,-1) -- (17,-2);
\node at (16.7,-1.25) {\tiny \textbf{F}};
\node at (16.28,-1.73) {\tiny \textbf{$\overline{\mathrm{F}}$}};
\node at (16.18,-1.18) {\tiny 9};
\node[fill=white, draw=black, circle, inner sep=0.4pt] at (16.8,-1.8) {\tiny \textbf{?}};
\fill[green!18, rounded corners=0.7pt] (17,-1) -- (18,-1) -- (18,-2) -- cycle;
\fill[gray!8, rounded corners=0.7pt] (17,-1) -- (17,-2) -- (18,-2) -- cycle;
\draw[black, rounded corners=0.7pt] (17,-1) rectangle (18,-2);
\draw[black, dashed, very thin] (17,-1) -- (18,-2);
\node at (17.7,-1.25) {\tiny \textbf{Ne}};
\node at (17.28,-1.73) {\tiny \textbf{$\overline{\mathrm{Ne}}$}};
\node at (17.18,-1.18) {\tiny 10};
\node[fill=white, draw=black, circle, inner sep=0.4pt] at (17.8,-1.8) {\tiny \textbf{?}};
\fill[blue!18, rounded corners=0.7pt] (0,-2) -- (1,-2) -- (1,-3) -- cycle;
\fill[gray!8, rounded corners=0.7pt] (0,-2) -- (0,-3) -- (1,-3) -- cycle;
\draw[black, rounded corners=0.7pt] (0,-2) rectangle (1,-3);
\draw[black, dashed, very thin] (0,-2) -- (1,-3);
\node at (0.7,-2.25) {\tiny \textbf{Na}};
\node at (0.28,-2.73) {\tiny \textbf{$\overline{\mathrm{Na}}$}};
\node at (0.18,-2.18) {\tiny 11};
\node[fill=white, draw=black, circle, inner sep=0.4pt] at (0.8,-2.8) {\tiny \textbf{?}};
\fill[blue!18, rounded corners=0.7pt] (1,-2) -- (2,-2) -- (2,-3) -- cycle;
\fill[gray!8, rounded corners=0.7pt] (1,-2) -- (1,-3) -- (2,-3) -- cycle;
\draw[black, rounded corners=0.7pt] (1,-2) rectangle (2,-3);
\draw[black, dashed, very thin] (1,-2) -- (2,-3);
\node at (1.7,-2.25) {\tiny \textbf{Mg}};
\node at (1.28,-2.73) {\tiny \textbf{$\overline{\mathrm{Mg}}$}};
\node at (1.18,-2.18) {\tiny 12};
\node[fill=white, draw=black, circle, inner sep=0.4pt] at (1.8,-2.8) {\tiny \textbf{?}};
\fill[green!18, rounded corners=0.7pt] (12,-2) -- (13,-2) -- (13,-3) -- cycle;
\fill[gray!8, rounded corners=0.7pt] (12,-2) -- (12,-3) -- (13,-3) -- cycle;
\draw[black, rounded corners=0.7pt] (12,-2) rectangle (13,-3);
\draw[black, dashed, very thin] (12,-2) -- (13,-3);
\node at (12.7,-2.25) {\tiny \textbf{Al}};
\node at (12.28,-2.73) {\tiny \textbf{$\overline{\mathrm{Al}}$}};
\node at (12.18,-2.18) {\tiny 13};
\node[fill=white, draw=black, circle, inner sep=0.4pt] at (12.8,-2.8) {\tiny \textbf{?}};
\fill[green!18, rounded corners=0.7pt] (13,-2) -- (14,-2) -- (14,-3) -- cycle;
\fill[gray!8, rounded corners=0.7pt] (13,-2) -- (13,-3) -- (14,-3) -- cycle;
\draw[black, rounded corners=0.7pt] (13,-2) rectangle (14,-3);
\draw[black, dashed, very thin] (13,-2) -- (14,-3);
\node at (13.7,-2.25) {\tiny \textbf{Si}};
\node at (13.28,-2.73) {\tiny \textbf{$\overline{\mathrm{Si}}$}};
\node at (13.18,-2.18) {\tiny 14};
\node[fill=white, draw=black, circle, inner sep=0.4pt] at (13.8,-2.8) {\tiny \textbf{?}};
\fill[green!18, rounded corners=0.7pt] (14,-2) -- (15,-2) -- (15,-3) -- cycle;
\fill[gray!8, rounded corners=0.7pt] (14,-2) -- (14,-3) -- (15,-3) -- cycle;
\draw[black, rounded corners=0.7pt] (14,-2) rectangle (15,-3);
\draw[black, dashed, very thin] (14,-2) -- (15,-3);
\node at (14.7,-2.25) {\tiny \textbf{P}};
\node at (14.28,-2.73) {\tiny \textbf{$\overline{\mathrm{P}}$}};
\node at (14.18,-2.18) {\tiny 15};
\node[fill=white, draw=black, circle, inner sep=0.4pt] at (14.8,-2.8) {\tiny \textbf{?}};
\fill[green!18, rounded corners=0.7pt] (15,-2) -- (16,-2) -- (16,-3) -- cycle;
\fill[gray!8, rounded corners=0.7pt] (15,-2) -- (15,-3) -- (16,-3) -- cycle;
\draw[black, rounded corners=0.7pt] (15,-2) rectangle (16,-3);
\draw[black, dashed, very thin] (15,-2) -- (16,-3);
\node at (15.7,-2.25) {\tiny \textbf{S}};
\node at (15.28,-2.73) {\tiny \textbf{$\overline{\mathrm{S}}$}};
\node at (15.18,-2.18) {\tiny 16};
\node[fill=white, draw=black, circle, inner sep=0.4pt] at (15.8,-2.8) {\tiny \textbf{?}};
\fill[green!18, rounded corners=0.7pt] (16,-2) -- (17,-2) -- (17,-3) -- cycle;
\fill[gray!8, rounded corners=0.7pt] (16,-2) -- (16,-3) -- (17,-3) -- cycle;
\draw[black, rounded corners=0.7pt] (16,-2) rectangle (17,-3);
\draw[black, dashed, very thin] (16,-2) -- (17,-3);
\node at (16.7,-2.25) {\tiny \textbf{Cl}};
\node at (16.28,-2.73) {\tiny \textbf{$\overline{\mathrm{Cl}}$}};
\node at (16.18,-2.18) {\tiny 17};
\node[fill=white, draw=black, circle, inner sep=0.4pt] at (16.8,-2.8) {\tiny \textbf{?}};
\fill[green!18, rounded corners=0.7pt] (17,-2) -- (18,-2) -- (18,-3) -- cycle;
\fill[gray!8, rounded corners=0.7pt] (17,-2) -- (17,-3) -- (18,-3) -- cycle;
\draw[black, rounded corners=0.7pt] (17,-2) rectangle (18,-3);
\draw[black, dashed, very thin] (17,-2) -- (18,-3);
\node at (17.7,-2.25) {\tiny \textbf{Ar}};
\node at (17.28,-2.73) {\tiny \textbf{$\overline{\mathrm{Ar}}$}};
\node at (17.18,-2.18) {\tiny 18};
\node[fill=white, draw=black, circle, inner sep=0.4pt] at (17.8,-2.8) {\tiny \textbf{?}};
\fill[blue!18, rounded corners=0.7pt] (0,-3) -- (1,-3) -- (1,-4) -- cycle;
\fill[gray!8, rounded corners=0.7pt] (0,-3) -- (0,-4) -- (1,-4) -- cycle;
\draw[black, rounded corners=0.7pt] (0,-3) rectangle (1,-4);
\draw[black, dashed, very thin] (0,-3) -- (1,-4);
\node at (0.7,-3.25) {\tiny \textbf{K}};
\node at (0.28,-3.73) {\tiny \textbf{$\overline{\mathrm{K}}$}};
\node at (0.18,-3.18) {\tiny 19};
\node[fill=white, draw=black, circle, inner sep=0.4pt] at (0.8,-3.8) {\tiny \textbf{?}};
\fill[blue!18, rounded corners=0.7pt] (1,-3) -- (2,-3) -- (2,-4) -- cycle;
\fill[gray!8, rounded corners=0.7pt] (1,-3) -- (1,-4) -- (2,-4) -- cycle;
\draw[black, rounded corners=0.7pt] (1,-3) rectangle (2,-4);
\draw[black, dashed, very thin] (1,-3) -- (2,-4);
\node at (1.7,-3.25) {\tiny \textbf{Ca}};
\node at (1.28,-3.73) {\tiny \textbf{$\overline{\mathrm{Ca}}$}};
\node at (1.18,-3.18) {\tiny 20};
\node[fill=white, draw=black, circle, inner sep=0.4pt] at (1.8,-3.8) {\tiny \textbf{?}};
\fill[orange!20, rounded corners=0.7pt] (2,-3) -- (3,-3) -- (3,-4) -- cycle;
\fill[gray!8, rounded corners=0.7pt] (2,-3) -- (2,-4) -- (3,-4) -- cycle;
\draw[black, rounded corners=0.7pt] (2,-3) rectangle (3,-4);
\draw[black, dashed, very thin] (2,-3) -- (3,-4);
\node at (2.7,-3.25) {\tiny \textbf{Sc}};
\node at (2.2800000000000002,-3.73) {\tiny \textbf{$\overline{\mathrm{Sc}}$}};
\node at (2.18,-3.18) {\tiny 21};
\node[fill=white, draw=black, circle, inner sep=0.4pt] at (2.8,-3.8) {\tiny \textbf{?}};
\fill[orange!20, rounded corners=0.7pt] (3,-3) -- (4,-3) -- (4,-4) -- cycle;
\fill[gray!8, rounded corners=0.7pt] (3,-3) -- (3,-4) -- (4,-4) -- cycle;
\draw[black, rounded corners=0.7pt] (3,-3) rectangle (4,-4);
\draw[black, dashed, very thin] (3,-3) -- (4,-4);
\node at (3.7,-3.25) {\tiny \textbf{Ti}};
\node at (3.2800000000000002,-3.73) {\tiny \textbf{$\overline{\mathrm{Ti}}$}};
\node at (3.18,-3.18) {\tiny 22};
\node[fill=white, draw=black, circle, inner sep=0.4pt] at (3.8,-3.8) {\tiny \textbf{?}};
\fill[orange!20, rounded corners=0.7pt] (4,-3) -- (5,-3) -- (5,-4) -- cycle;
\fill[gray!8, rounded corners=0.7pt] (4,-3) -- (4,-4) -- (5,-4) -- cycle;
\draw[black, rounded corners=0.7pt] (4,-3) rectangle (5,-4);
\draw[black, dashed, very thin] (4,-3) -- (5,-4);
\node at (4.7,-3.25) {\tiny \textbf{V}};
\node at (4.28,-3.73) {\tiny \textbf{$\overline{\mathrm{V}}$}};
\node at (4.18,-3.18) {\tiny 23};
\node[fill=white, draw=black, circle, inner sep=0.4pt] at (4.8,-3.8) {\tiny \textbf{?}};
\fill[orange!20, rounded corners=0.7pt] (5,-3) -- (6,-3) -- (6,-4) -- cycle;
\fill[gray!8, rounded corners=0.7pt] (5,-3) -- (5,-4) -- (6,-4) -- cycle;
\draw[black, rounded corners=0.7pt] (5,-3) rectangle (6,-4);
\draw[black, dashed, very thin] (5,-3) -- (6,-4);
\node at (5.7,-3.25) {\tiny \textbf{Cr}};
\node at (5.28,-3.73) {\tiny \textbf{$\overline{\mathrm{Cr}}$}};
\node at (5.18,-3.18) {\tiny 24};
\node[fill=white, draw=black, circle, inner sep=0.4pt] at (5.8,-3.8) {\tiny \textbf{?}};
\fill[orange!20, rounded corners=0.7pt] (6,-3) -- (7,-3) -- (7,-4) -- cycle;
\fill[gray!8, rounded corners=0.7pt] (6,-3) -- (6,-4) -- (7,-4) -- cycle;
\draw[black, rounded corners=0.7pt] (6,-3) rectangle (7,-4);
\draw[black, dashed, very thin] (6,-3) -- (7,-4);
\node at (6.7,-3.25) {\tiny \textbf{Mn}};
\node at (6.28,-3.73) {\tiny \textbf{$\overline{\mathrm{Mn}}$}};
\node at (6.18,-3.18) {\tiny 25};
\node[fill=white, draw=black, circle, inner sep=0.4pt] at (6.8,-3.8) {\tiny \textbf{?}};
\fill[orange!20, rounded corners=0.7pt] (7,-3) -- (8,-3) -- (8,-4) -- cycle;
\fill[gray!8, rounded corners=0.7pt] (7,-3) -- (7,-4) -- (8,-4) -- cycle;
\draw[black, rounded corners=0.7pt] (7,-3) rectangle (8,-4);
\draw[black, dashed, very thin] (7,-3) -- (8,-4);
\node at (7.7,-3.25) {\tiny \textbf{Fe}};
\node at (7.28,-3.73) {\tiny \textbf{$\overline{\mathrm{Fe}}$}};
\node at (7.18,-3.18) {\tiny 26};
\node[fill=white, draw=black, circle, inner sep=0.4pt] at (7.8,-3.8) {\tiny \textbf{?}};
\fill[orange!20, rounded corners=0.7pt] (8,-3) -- (9,-3) -- (9,-4) -- cycle;
\fill[gray!8, rounded corners=0.7pt] (8,-3) -- (8,-4) -- (9,-4) -- cycle;
\draw[black, rounded corners=0.7pt] (8,-3) rectangle (9,-4);
\draw[black, dashed, very thin] (8,-3) -- (9,-4);
\node at (8.7,-3.25) {\tiny \textbf{Co}};
\node at (8.28,-3.73) {\tiny \textbf{$\overline{\mathrm{Co}}$}};
\node at (8.18,-3.18) {\tiny 27};
\node[fill=white, draw=black, circle, inner sep=0.4pt] at (8.8,-3.8) {\tiny \textbf{?}};
\fill[orange!20, rounded corners=0.7pt] (9,-3) -- (10,-3) -- (10,-4) -- cycle;
\fill[gray!8, rounded corners=0.7pt] (9,-3) -- (9,-4) -- (10,-4) -- cycle;
\draw[black, rounded corners=0.7pt] (9,-3) rectangle (10,-4);
\draw[black, dashed, very thin] (9,-3) -- (10,-4);
\node at (9.7,-3.25) {\tiny \textbf{Ni}};
\node at (9.28,-3.73) {\tiny \textbf{$\overline{\mathrm{Ni}}$}};
\node at (9.18,-3.18) {\tiny 28};
\node[fill=white, draw=black, circle, inner sep=0.4pt] at (9.8,-3.8) {\tiny \textbf{?}};
\fill[orange!20, rounded corners=0.7pt] (10,-3) -- (11,-3) -- (11,-4) -- cycle;
\fill[gray!8, rounded corners=0.7pt] (10,-3) -- (10,-4) -- (11,-4) -- cycle;
\draw[black, rounded corners=0.7pt] (10,-3) rectangle (11,-4);
\draw[black, dashed, very thin] (10,-3) -- (11,-4);
\node at (10.7,-3.25) {\tiny \textbf{Cu}};
\node at (10.28,-3.73) {\tiny \textbf{$\overline{\mathrm{Cu}}$}};
\node at (10.18,-3.18) {\tiny 29};
\node[fill=white, draw=black, circle, inner sep=0.4pt] at (10.8,-3.8) {\tiny \textbf{?}};
\fill[orange!20, rounded corners=0.7pt] (11,-3) -- (12,-3) -- (12,-4) -- cycle;
\fill[gray!8, rounded corners=0.7pt] (11,-3) -- (11,-4) -- (12,-4) -- cycle;
\draw[black, rounded corners=0.7pt] (11,-3) rectangle (12,-4);
\draw[black, dashed, very thin] (11,-3) -- (12,-4);
\node at (11.7,-3.25) {\tiny \textbf{Zn}};
\node at (11.28,-3.73) {\tiny \textbf{$\overline{\mathrm{Zn}}$}};
\node at (11.18,-3.18) {\tiny 30};
\node[fill=white, draw=black, circle, inner sep=0.4pt] at (11.8,-3.8) {\tiny \textbf{?}};
\fill[green!18, rounded corners=0.7pt] (12,-3) -- (13,-3) -- (13,-4) -- cycle;
\fill[gray!8, rounded corners=0.7pt] (12,-3) -- (12,-4) -- (13,-4) -- cycle;
\draw[black, rounded corners=0.7pt] (12,-3) rectangle (13,-4);
\draw[black, dashed, very thin] (12,-3) -- (13,-4);
\node at (12.7,-3.25) {\tiny \textbf{Ga}};
\node at (12.28,-3.73) {\tiny \textbf{$\overline{\mathrm{Ga}}$}};
\node at (12.18,-3.18) {\tiny 31};
\node[fill=white, draw=black, circle, inner sep=0.4pt] at (12.8,-3.8) {\tiny \textbf{?}};
\fill[green!18, rounded corners=0.7pt] (13,-3) -- (14,-3) -- (14,-4) -- cycle;
\fill[gray!8, rounded corners=0.7pt] (13,-3) -- (13,-4) -- (14,-4) -- cycle;
\draw[black, rounded corners=0.7pt] (13,-3) rectangle (14,-4);
\draw[black, dashed, very thin] (13,-3) -- (14,-4);
\node at (13.7,-3.25) {\tiny \textbf{Ge}};
\node at (13.28,-3.73) {\tiny \textbf{$\overline{\mathrm{Ge}}$}};
\node at (13.18,-3.18) {\tiny 32};
\node[fill=white, draw=black, circle, inner sep=0.4pt] at (13.8,-3.8) {\tiny \textbf{?}};
\fill[green!18, rounded corners=0.7pt] (14,-3) -- (15,-3) -- (15,-4) -- cycle;
\fill[gray!8, rounded corners=0.7pt] (14,-3) -- (14,-4) -- (15,-4) -- cycle;
\draw[black, rounded corners=0.7pt] (14,-3) rectangle (15,-4);
\draw[black, dashed, very thin] (14,-3) -- (15,-4);
\node at (14.7,-3.25) {\tiny \textbf{As}};
\node at (14.28,-3.73) {\tiny \textbf{$\overline{\mathrm{As}}$}};
\node at (14.18,-3.18) {\tiny 33};
\node[fill=white, draw=black, circle, inner sep=0.4pt] at (14.8,-3.8) {\tiny \textbf{?}};
\fill[green!18, rounded corners=0.7pt] (15,-3) -- (16,-3) -- (16,-4) -- cycle;
\fill[gray!8, rounded corners=0.7pt] (15,-3) -- (15,-4) -- (16,-4) -- cycle;
\draw[black, rounded corners=0.7pt] (15,-3) rectangle (16,-4);
\draw[black, dashed, very thin] (15,-3) -- (16,-4);
\node at (15.7,-3.25) {\tiny \textbf{Se}};
\node at (15.28,-3.73) {\tiny \textbf{$\overline{\mathrm{Se}}$}};
\node at (15.18,-3.18) {\tiny 34};
\node[fill=white, draw=black, circle, inner sep=0.4pt] at (15.8,-3.8) {\tiny \textbf{?}};
\fill[green!18, rounded corners=0.7pt] (16,-3) -- (17,-3) -- (17,-4) -- cycle;
\fill[gray!8, rounded corners=0.7pt] (16,-3) -- (16,-4) -- (17,-4) -- cycle;
\draw[black, rounded corners=0.7pt] (16,-3) rectangle (17,-4);
\draw[black, dashed, very thin] (16,-3) -- (17,-4);
\node at (16.7,-3.25) {\tiny \textbf{Br}};
\node at (16.28,-3.73) {\tiny \textbf{$\overline{\mathrm{Br}}$}};
\node at (16.18,-3.18) {\tiny 35};
\node[fill=white, draw=black, circle, inner sep=0.4pt] at (16.8,-3.8) {\tiny \textbf{?}};
\fill[green!18, rounded corners=0.7pt] (17,-3) -- (18,-3) -- (18,-4) -- cycle;
\fill[gray!8, rounded corners=0.7pt] (17,-3) -- (17,-4) -- (18,-4) -- cycle;
\draw[black, rounded corners=0.7pt] (17,-3) rectangle (18,-4);
\draw[black, dashed, very thin] (17,-3) -- (18,-4);
\node at (17.7,-3.25) {\tiny \textbf{Kr}};
\node at (17.28,-3.73) {\tiny \textbf{$\overline{\mathrm{Kr}}$}};
\node at (17.18,-3.18) {\tiny 36};
\node[fill=white, draw=black, circle, inner sep=0.4pt] at (17.8,-3.8) {\tiny \textbf{?}};
\fill[blue!18, rounded corners=0.7pt] (0,-4) -- (1,-4) -- (1,-5) -- cycle;
\fill[gray!8, rounded corners=0.7pt] (0,-4) -- (0,-5) -- (1,-5) -- cycle;
\draw[black, rounded corners=0.7pt] (0,-4) rectangle (1,-5);
\draw[black, dashed, very thin] (0,-4) -- (1,-5);
\node at (0.7,-4.25) {\tiny \textbf{Rb}};
\node at (0.28,-4.73) {\tiny \textbf{$\overline{\mathrm{Rb}}$}};
\node at (0.18,-4.18) {\tiny 37};
\node[fill=white, draw=black, circle, inner sep=0.4pt] at (0.8,-4.8) {\tiny \textbf{?}};
\fill[blue!18, rounded corners=0.7pt] (1,-4) -- (2,-4) -- (2,-5) -- cycle;
\fill[gray!8, rounded corners=0.7pt] (1,-4) -- (1,-5) -- (2,-5) -- cycle;
\draw[black, rounded corners=0.7pt] (1,-4) rectangle (2,-5);
\draw[black, dashed, very thin] (1,-4) -- (2,-5);
\node at (1.7,-4.25) {\tiny \textbf{Sr}};
\node at (1.28,-4.73) {\tiny \textbf{$\overline{\mathrm{Sr}}$}};
\node at (1.18,-4.18) {\tiny 38};
\node[fill=white, draw=black, circle, inner sep=0.4pt] at (1.8,-4.8) {\tiny \textbf{?}};
\fill[orange!20, rounded corners=0.7pt] (2,-4) -- (3,-4) -- (3,-5) -- cycle;
\fill[gray!8, rounded corners=0.7pt] (2,-4) -- (2,-5) -- (3,-5) -- cycle;
\draw[black, rounded corners=0.7pt] (2,-4) rectangle (3,-5);
\draw[black, dashed, very thin] (2,-4) -- (3,-5);
\node at (2.7,-4.25) {\tiny \textbf{Y}};
\node at (2.2800000000000002,-4.73) {\tiny \textbf{$\overline{\mathrm{Y}}$}};
\node at (2.18,-4.18) {\tiny 39};
\node[fill=white, draw=black, circle, inner sep=0.4pt] at (2.8,-4.8) {\tiny \textbf{?}};
\fill[orange!20, rounded corners=0.7pt] (3,-4) -- (4,-4) -- (4,-5) -- cycle;
\fill[gray!8, rounded corners=0.7pt] (3,-4) -- (3,-5) -- (4,-5) -- cycle;
\draw[black, rounded corners=0.7pt] (3,-4) rectangle (4,-5);
\draw[black, dashed, very thin] (3,-4) -- (4,-5);
\node at (3.7,-4.25) {\tiny \textbf{Zr}};
\node at (3.2800000000000002,-4.73) {\tiny \textbf{$\overline{\mathrm{Zr}}$}};
\node at (3.18,-4.18) {\tiny 40};
\node[fill=white, draw=black, circle, inner sep=0.4pt] at (3.8,-4.8) {\tiny \textbf{?}};
\fill[orange!20, rounded corners=0.7pt] (4,-4) -- (5,-4) -- (5,-5) -- cycle;
\fill[gray!8, rounded corners=0.7pt] (4,-4) -- (4,-5) -- (5,-5) -- cycle;
\draw[black, rounded corners=0.7pt] (4,-4) rectangle (5,-5);
\draw[black, dashed, very thin] (4,-4) -- (5,-5);
\node at (4.7,-4.25) {\tiny \textbf{Nb}};
\node at (4.28,-4.73) {\tiny \textbf{$\overline{\mathrm{Nb}}$}};
\node at (4.18,-4.18) {\tiny 41};
\node[fill=white, draw=black, circle, inner sep=0.4pt] at (4.8,-4.8) {\tiny \textbf{?}};
\fill[orange!20, rounded corners=0.7pt] (5,-4) -- (6,-4) -- (6,-5) -- cycle;
\fill[gray!8, rounded corners=0.7pt] (5,-4) -- (5,-5) -- (6,-5) -- cycle;
\draw[black, rounded corners=0.7pt] (5,-4) rectangle (6,-5);
\draw[black, dashed, very thin] (5,-4) -- (6,-5);
\node at (5.7,-4.25) {\tiny \textbf{Mo}};
\node at (5.28,-4.73) {\tiny \textbf{$\overline{\mathrm{Mo}}$}};
\node at (5.18,-4.18) {\tiny 42};
\node[fill=white, draw=black, circle, inner sep=0.4pt] at (5.8,-4.8) {\tiny \textbf{?}};
\fill[orange!20, rounded corners=0.7pt] (6,-4) -- (7,-4) -- (7,-5) -- cycle;
\fill[gray!8, rounded corners=0.7pt] (6,-4) -- (6,-5) -- (7,-5) -- cycle;
\draw[black, rounded corners=0.7pt] (6,-4) rectangle (7,-5);
\draw[black, dashed, very thin] (6,-4) -- (7,-5);
\node at (6.7,-4.25) {\tiny \textbf{Tc}};
\node at (6.28,-4.73) {\tiny \textbf{$\overline{\mathrm{Tc}}$}};
\node at (6.18,-4.18) {\tiny 43};
\node[fill=white, draw=black, circle, inner sep=0.4pt] at (6.8,-4.8) {\tiny \textbf{?}};
\fill[orange!20, rounded corners=0.7pt] (7,-4) -- (8,-4) -- (8,-5) -- cycle;
\fill[gray!8, rounded corners=0.7pt] (7,-4) -- (7,-5) -- (8,-5) -- cycle;
\draw[black, rounded corners=0.7pt] (7,-4) rectangle (8,-5);
\draw[black, dashed, very thin] (7,-4) -- (8,-5);
\node at (7.7,-4.25) {\tiny \textbf{Ru}};
\node at (7.28,-4.73) {\tiny \textbf{$\overline{\mathrm{Ru}}$}};
\node at (7.18,-4.18) {\tiny 44};
\node[fill=white, draw=black, circle, inner sep=0.4pt] at (7.8,-4.8) {\tiny \textbf{?}};
\fill[orange!20, rounded corners=0.7pt] (8,-4) -- (9,-4) -- (9,-5) -- cycle;
\fill[gray!8, rounded corners=0.7pt] (8,-4) -- (8,-5) -- (9,-5) -- cycle;
\draw[black, rounded corners=0.7pt] (8,-4) rectangle (9,-5);
\draw[black, dashed, very thin] (8,-4) -- (9,-5);
\node at (8.7,-4.25) {\tiny \textbf{Rh}};
\node at (8.28,-4.73) {\tiny \textbf{$\overline{\mathrm{Rh}}$}};
\node at (8.18,-4.18) {\tiny 45};
\node[fill=white, draw=black, circle, inner sep=0.4pt] at (8.8,-4.8) {\tiny \textbf{?}};
\fill[orange!20, rounded corners=0.7pt] (9,-4) -- (10,-4) -- (10,-5) -- cycle;
\fill[gray!8, rounded corners=0.7pt] (9,-4) -- (9,-5) -- (10,-5) -- cycle;
\draw[black, rounded corners=0.7pt] (9,-4) rectangle (10,-5);
\draw[black, dashed, very thin] (9,-4) -- (10,-5);
\node at (9.7,-4.25) {\tiny \textbf{Pd}};
\node at (9.28,-4.73) {\tiny \textbf{$\overline{\mathrm{Pd}}$}};
\node at (9.18,-4.18) {\tiny 46};
\node[fill=white, draw=black, circle, inner sep=0.4pt] at (9.8,-4.8) {\tiny \textbf{?}};
\fill[orange!20, rounded corners=0.7pt] (10,-4) -- (11,-4) -- (11,-5) -- cycle;
\fill[gray!8, rounded corners=0.7pt] (10,-4) -- (10,-5) -- (11,-5) -- cycle;
\draw[black, rounded corners=0.7pt] (10,-4) rectangle (11,-5);
\draw[black, dashed, very thin] (10,-4) -- (11,-5);
\node at (10.7,-4.25) {\tiny \textbf{Ag}};
\node at (10.28,-4.73) {\tiny \textbf{$\overline{\mathrm{Ag}}$}};
\node at (10.18,-4.18) {\tiny 47};
\node[fill=white, draw=black, circle, inner sep=0.4pt] at (10.8,-4.8) {\tiny \textbf{?}};
\fill[orange!20, rounded corners=0.7pt] (11,-4) -- (12,-4) -- (12,-5) -- cycle;
\fill[gray!8, rounded corners=0.7pt] (11,-4) -- (11,-5) -- (12,-5) -- cycle;
\draw[black, rounded corners=0.7pt] (11,-4) rectangle (12,-5);
\draw[black, dashed, very thin] (11,-4) -- (12,-5);
\node at (11.7,-4.25) {\tiny \textbf{Cd}};
\node at (11.28,-4.73) {\tiny \textbf{$\overline{\mathrm{Cd}}$}};
\node at (11.18,-4.18) {\tiny 48};
\node[fill=white, draw=black, circle, inner sep=0.4pt] at (11.8,-4.8) {\tiny \textbf{?}};
\fill[green!18, rounded corners=0.7pt] (12,-4) -- (13,-4) -- (13,-5) -- cycle;
\fill[gray!8, rounded corners=0.7pt] (12,-4) -- (12,-5) -- (13,-5) -- cycle;
\draw[black, rounded corners=0.7pt] (12,-4) rectangle (13,-5);
\draw[black, dashed, very thin] (12,-4) -- (13,-5);
\node at (12.7,-4.25) {\tiny \textbf{In}};
\node at (12.28,-4.73) {\tiny \textbf{$\overline{\mathrm{In}}$}};
\node at (12.18,-4.18) {\tiny 49};
\node[fill=white, draw=black, circle, inner sep=0.4pt] at (12.8,-4.8) {\tiny \textbf{?}};
\fill[green!18, rounded corners=0.7pt] (13,-4) -- (14,-4) -- (14,-5) -- cycle;
\fill[gray!8, rounded corners=0.7pt] (13,-4) -- (13,-5) -- (14,-5) -- cycle;
\draw[black, rounded corners=0.7pt] (13,-4) rectangle (14,-5);
\draw[black, dashed, very thin] (13,-4) -- (14,-5);
\node at (13.7,-4.25) {\tiny \textbf{Sn}};
\node at (13.28,-4.73) {\tiny \textbf{$\overline{\mathrm{Sn}}$}};
\node at (13.18,-4.18) {\tiny 50};
\node[fill=white, draw=black, circle, inner sep=0.4pt] at (13.8,-4.8) {\tiny \textbf{?}};
\fill[green!18, rounded corners=0.7pt] (14,-4) -- (15,-4) -- (15,-5) -- cycle;
\fill[gray!8, rounded corners=0.7pt] (14,-4) -- (14,-5) -- (15,-5) -- cycle;
\draw[black, rounded corners=0.7pt] (14,-4) rectangle (15,-5);
\draw[black, dashed, very thin] (14,-4) -- (15,-5);
\node at (14.7,-4.25) {\tiny \textbf{Sb}};
\node at (14.28,-4.73) {\tiny \textbf{$\overline{\mathrm{Sb}}$}};
\node at (14.18,-4.18) {\tiny 51};
\node[fill=white, draw=black, circle, inner sep=0.4pt] at (14.8,-4.8) {\tiny \textbf{?}};
\fill[green!18, rounded corners=0.7pt] (15,-4) -- (16,-4) -- (16,-5) -- cycle;
\fill[gray!8, rounded corners=0.7pt] (15,-4) -- (15,-5) -- (16,-5) -- cycle;
\draw[black, rounded corners=0.7pt] (15,-4) rectangle (16,-5);
\draw[black, dashed, very thin] (15,-4) -- (16,-5);
\node at (15.7,-4.25) {\tiny \textbf{Te}};
\node at (15.28,-4.73) {\tiny \textbf{$\overline{\mathrm{Te}}$}};
\node at (15.18,-4.18) {\tiny 52};
\node[fill=white, draw=black, circle, inner sep=0.4pt] at (15.8,-4.8) {\tiny \textbf{?}};
\fill[green!18, rounded corners=0.7pt] (16,-4) -- (17,-4) -- (17,-5) -- cycle;
\fill[gray!8, rounded corners=0.7pt] (16,-4) -- (16,-5) -- (17,-5) -- cycle;
\draw[black, rounded corners=0.7pt] (16,-4) rectangle (17,-5);
\draw[black, dashed, very thin] (16,-4) -- (17,-5);
\node at (16.7,-4.25) {\tiny \textbf{I}};
\node at (16.28,-4.73) {\tiny \textbf{$\overline{\mathrm{I}}$}};
\node at (16.18,-4.18) {\tiny 53};
\node[fill=white, draw=black, circle, inner sep=0.4pt] at (16.8,-4.8) {\tiny \textbf{?}};
\fill[green!18, rounded corners=0.7pt] (17,-4) -- (18,-4) -- (18,-5) -- cycle;
\fill[gray!8, rounded corners=0.7pt] (17,-4) -- (17,-5) -- (18,-5) -- cycle;
\draw[black, rounded corners=0.7pt] (17,-4) rectangle (18,-5);
\draw[black, dashed, very thin] (17,-4) -- (18,-5);
\node at (17.7,-4.25) {\tiny \textbf{Xe}};
\node at (17.28,-4.73) {\tiny \textbf{$\overline{\mathrm{Xe}}$}};
\node at (17.18,-4.18) {\tiny 54};
\node[fill=white, draw=black, circle, inner sep=0.4pt] at (17.8,-4.8) {\tiny \textbf{?}};
\fill[blue!18, rounded corners=0.7pt] (0,-5) -- (1,-5) -- (1,-6) -- cycle;
\fill[gray!8, rounded corners=0.7pt] (0,-5) -- (0,-6) -- (1,-6) -- cycle;
\draw[black, rounded corners=0.7pt] (0,-5) rectangle (1,-6);
\draw[black, dashed, very thin] (0,-5) -- (1,-6);
\node at (0.7,-5.25) {\tiny \textbf{Cs}};
\node at (0.28,-5.73) {\tiny \textbf{$\overline{\mathrm{Cs}}$}};
\node at (0.18,-5.18) {\tiny 55};
\node[fill=white, draw=black, circle, inner sep=0.4pt] at (0.8,-5.8) {\tiny \textbf{?}};
\fill[blue!18, rounded corners=0.7pt] (1,-5) -- (2,-5) -- (2,-6) -- cycle;
\fill[gray!8, rounded corners=0.7pt] (1,-5) -- (1,-6) -- (2,-6) -- cycle;
\draw[black, rounded corners=0.7pt] (1,-5) rectangle (2,-6);
\draw[black, dashed, very thin] (1,-5) -- (2,-6);
\node at (1.7,-5.25) {\tiny \textbf{Ba}};
\node at (1.28,-5.73) {\tiny \textbf{$\overline{\mathrm{Ba}}$}};
\node at (1.18,-5.18) {\tiny 56};
\node[fill=white, draw=black, circle, inner sep=0.4pt] at (1.8,-5.8) {\tiny \textbf{?}};
\fill[orange!20, rounded corners=0.7pt] (2,-5) -- (3,-5) -- (3,-6) -- cycle;
\fill[gray!8, rounded corners=0.7pt] (2,-5) -- (2,-6) -- (3,-6) -- cycle;
\draw[black, rounded corners=0.7pt] (2,-5) rectangle (3,-6);
\draw[black, dashed, very thin] (2,-5) -- (3,-6);
\node at (2.7,-5.25) {\tiny \textbf{Lu}};
\node at (2.2800000000000002,-5.73) {\tiny \textbf{$\overline{\mathrm{Lu}}$}};
\node at (2.18,-5.18) {\tiny 71};
\node[fill=white, draw=black, circle, inner sep=0.4pt] at (2.8,-5.8) {\tiny \textbf{?}};
\fill[orange!20, rounded corners=0.7pt] (3,-5) -- (4,-5) -- (4,-6) -- cycle;
\fill[gray!8, rounded corners=0.7pt] (3,-5) -- (3,-6) -- (4,-6) -- cycle;
\draw[black, rounded corners=0.7pt] (3,-5) rectangle (4,-6);
\draw[black, dashed, very thin] (3,-5) -- (4,-6);
\node at (3.7,-5.25) {\tiny \textbf{Hf}};
\node at (3.2800000000000002,-5.73) {\tiny \textbf{$\overline{\mathrm{Hf}}$}};
\node at (3.18,-5.18) {\tiny 72};
\node[fill=white, draw=black, circle, inner sep=0.4pt] at (3.8,-5.8) {\tiny \textbf{?}};
\fill[orange!20, rounded corners=0.7pt] (4,-5) -- (5,-5) -- (5,-6) -- cycle;
\fill[gray!8, rounded corners=0.7pt] (4,-5) -- (4,-6) -- (5,-6) -- cycle;
\draw[black, rounded corners=0.7pt] (4,-5) rectangle (5,-6);
\draw[black, dashed, very thin] (4,-5) -- (5,-6);
\node at (4.7,-5.25) {\tiny \textbf{Ta}};
\node at (4.28,-5.73) {\tiny \textbf{$\overline{\mathrm{Ta}}$}};
\node at (4.18,-5.18) {\tiny 73};
\node[fill=white, draw=black, circle, inner sep=0.4pt] at (4.8,-5.8) {\tiny \textbf{?}};
\fill[orange!20, rounded corners=0.7pt] (5,-5) -- (6,-5) -- (6,-6) -- cycle;
\fill[gray!8, rounded corners=0.7pt] (5,-5) -- (5,-6) -- (6,-6) -- cycle;
\draw[black, rounded corners=0.7pt] (5,-5) rectangle (6,-6);
\draw[black, dashed, very thin] (5,-5) -- (6,-6);
\node at (5.7,-5.25) {\tiny \textbf{W}};
\node at (5.28,-5.73) {\tiny \textbf{$\overline{\mathrm{W}}$}};
\node at (5.18,-5.18) {\tiny 74};
\node[fill=white, draw=black, circle, inner sep=0.4pt] at (5.8,-5.8) {\tiny \textbf{?}};
\fill[orange!20, rounded corners=0.7pt] (6,-5) -- (7,-5) -- (7,-6) -- cycle;
\fill[gray!8, rounded corners=0.7pt] (6,-5) -- (6,-6) -- (7,-6) -- cycle;
\draw[black, rounded corners=0.7pt] (6,-5) rectangle (7,-6);
\draw[black, dashed, very thin] (6,-5) -- (7,-6);
\node at (6.7,-5.25) {\tiny \textbf{Re}};
\node at (6.28,-5.73) {\tiny \textbf{$\overline{\mathrm{Re}}$}};
\node at (6.18,-5.18) {\tiny 75};
\node[fill=white, draw=black, circle, inner sep=0.4pt] at (6.8,-5.8) {\tiny \textbf{?}};
\fill[orange!20, rounded corners=0.7pt] (7,-5) -- (8,-5) -- (8,-6) -- cycle;
\fill[gray!8, rounded corners=0.7pt] (7,-5) -- (7,-6) -- (8,-6) -- cycle;
\draw[black, rounded corners=0.7pt] (7,-5) rectangle (8,-6);
\draw[black, dashed, very thin] (7,-5) -- (8,-6);
\node at (7.7,-5.25) {\tiny \textbf{Os}};
\node at (7.28,-5.73) {\tiny \textbf{$\overline{\mathrm{Os}}$}};
\node at (7.18,-5.18) {\tiny 76};
\node[fill=white, draw=black, circle, inner sep=0.4pt] at (7.8,-5.8) {\tiny \textbf{?}};
\fill[orange!20, rounded corners=0.7pt] (8,-5) -- (9,-5) -- (9,-6) -- cycle;
\fill[gray!8, rounded corners=0.7pt] (8,-5) -- (8,-6) -- (9,-6) -- cycle;
\draw[black, rounded corners=0.7pt] (8,-5) rectangle (9,-6);
\draw[black, dashed, very thin] (8,-5) -- (9,-6);
\node at (8.7,-5.25) {\tiny \textbf{Ir}};
\node at (8.28,-5.73) {\tiny \textbf{$\overline{\mathrm{Ir}}$}};
\node at (8.18,-5.18) {\tiny 77};
\node[fill=white, draw=black, circle, inner sep=0.4pt] at (8.8,-5.8) {\tiny \textbf{?}};
\fill[orange!20, rounded corners=0.7pt] (9,-5) -- (10,-5) -- (10,-6) -- cycle;
\fill[gray!8, rounded corners=0.7pt] (9,-5) -- (9,-6) -- (10,-6) -- cycle;
\draw[black, rounded corners=0.7pt] (9,-5) rectangle (10,-6);
\draw[black, dashed, very thin] (9,-5) -- (10,-6);
\node at (9.7,-5.25) {\tiny \textbf{Pt}};
\node at (9.28,-5.73) {\tiny \textbf{$\overline{\mathrm{Pt}}$}};
\node at (9.18,-5.18) {\tiny 78};
\node[fill=white, draw=black, circle, inner sep=0.4pt] at (9.8,-5.8) {\tiny \textbf{?}};
\fill[orange!20, rounded corners=0.7pt] (10,-5) -- (11,-5) -- (11,-6) -- cycle;
\fill[gray!8, rounded corners=0.7pt] (10,-5) -- (10,-6) -- (11,-6) -- cycle;
\draw[black, rounded corners=0.7pt] (10,-5) rectangle (11,-6);
\draw[black, dashed, very thin] (10,-5) -- (11,-6);
\node at (10.7,-5.25) {\tiny \textbf{Au}};
\node at (10.28,-5.73) {\tiny \textbf{$\overline{\mathrm{Au}}$}};
\node at (10.18,-5.18) {\tiny 79};
\node[fill=white, draw=black, circle, inner sep=0.4pt] at (10.8,-5.8) {\tiny \textbf{?}};
\fill[orange!20, rounded corners=0.7pt] (11,-5) -- (12,-5) -- (12,-6) -- cycle;
\fill[gray!8, rounded corners=0.7pt] (11,-5) -- (11,-6) -- (12,-6) -- cycle;
\draw[black, rounded corners=0.7pt] (11,-5) rectangle (12,-6);
\draw[black, dashed, very thin] (11,-5) -- (12,-6);
\node at (11.7,-5.25) {\tiny \textbf{Hg}};
\node at (11.28,-5.73) {\tiny \textbf{$\overline{\mathrm{Hg}}$}};
\node at (11.18,-5.18) {\tiny 80};
\node[fill=white, draw=black, circle, inner sep=0.4pt] at (11.8,-5.8) {\tiny \textbf{?}};
\fill[green!18, rounded corners=0.7pt] (12,-5) -- (13,-5) -- (13,-6) -- cycle;
\fill[gray!8, rounded corners=0.7pt] (12,-5) -- (12,-6) -- (13,-6) -- cycle;
\draw[black, rounded corners=0.7pt] (12,-5) rectangle (13,-6);
\draw[black, dashed, very thin] (12,-5) -- (13,-6);
\node at (12.7,-5.25) {\tiny \textbf{Tl}};
\node at (12.28,-5.73) {\tiny \textbf{$\overline{\mathrm{Tl}}$}};
\node at (12.18,-5.18) {\tiny 81};
\node[fill=white, draw=black, circle, inner sep=0.4pt] at (12.8,-5.8) {\tiny \textbf{?}};
\fill[green!18, rounded corners=0.7pt] (13,-5) -- (14,-5) -- (14,-6) -- cycle;
\fill[gray!8, rounded corners=0.7pt] (13,-5) -- (13,-6) -- (14,-6) -- cycle;
\draw[black, rounded corners=0.7pt] (13,-5) rectangle (14,-6);
\draw[black, dashed, very thin] (13,-5) -- (14,-6);
\node at (13.7,-5.25) {\tiny \textbf{Pb}};
\node at (13.28,-5.73) {\tiny \textbf{$\overline{\mathrm{Pb}}$}};
\node at (13.18,-5.18) {\tiny 82};
\node[fill=white, draw=black, circle, inner sep=0.4pt] at (13.8,-5.8) {\tiny \textbf{?}};
\fill[green!18, rounded corners=0.7pt] (14,-5) -- (15,-5) -- (15,-6) -- cycle;
\fill[gray!8, rounded corners=0.7pt] (14,-5) -- (14,-6) -- (15,-6) -- cycle;
\draw[black, rounded corners=0.7pt] (14,-5) rectangle (15,-6);
\draw[black, dashed, very thin] (14,-5) -- (15,-6);
\node at (14.7,-5.25) {\tiny \textbf{Bi}};
\node at (14.28,-5.73) {\tiny \textbf{$\overline{\mathrm{Bi}}$}};
\node at (14.18,-5.18) {\tiny 83};
\node[fill=white, draw=black, circle, inner sep=0.4pt] at (14.8,-5.8) {\tiny \textbf{?}};
\fill[green!18, rounded corners=0.7pt] (15,-5) -- (16,-5) -- (16,-6) -- cycle;
\fill[gray!8, rounded corners=0.7pt] (15,-5) -- (15,-6) -- (16,-6) -- cycle;
\draw[black, rounded corners=0.7pt] (15,-5) rectangle (16,-6);
\draw[black, dashed, very thin] (15,-5) -- (16,-6);
\node at (15.7,-5.25) {\tiny \textbf{Po}};
\node at (15.28,-5.73) {\tiny \textbf{$\overline{\mathrm{Po}}$}};
\node at (15.18,-5.18) {\tiny 84};
\node[fill=white, draw=black, circle, inner sep=0.4pt] at (15.8,-5.8) {\tiny \textbf{?}};
\fill[green!18, rounded corners=0.7pt] (16,-5) -- (17,-5) -- (17,-6) -- cycle;
\fill[gray!8, rounded corners=0.7pt] (16,-5) -- (16,-6) -- (17,-6) -- cycle;
\draw[black, rounded corners=0.7pt] (16,-5) rectangle (17,-6);
\draw[black, dashed, very thin] (16,-5) -- (17,-6);
\node at (16.7,-5.25) {\tiny \textbf{At}};
\node at (16.28,-5.73) {\tiny \textbf{$\overline{\mathrm{At}}$}};
\node at (16.18,-5.18) {\tiny 85};
\node[fill=white, draw=black, circle, inner sep=0.4pt] at (16.8,-5.8) {\tiny \textbf{?}};
\fill[green!18, rounded corners=0.7pt] (17,-5) -- (18,-5) -- (18,-6) -- cycle;
\fill[gray!8, rounded corners=0.7pt] (17,-5) -- (17,-6) -- (18,-6) -- cycle;
\draw[black, rounded corners=0.7pt] (17,-5) rectangle (18,-6);
\draw[black, dashed, very thin] (17,-5) -- (18,-6);
\node at (17.7,-5.25) {\tiny \textbf{Rn}};
\node at (17.28,-5.73) {\tiny \textbf{$\overline{\mathrm{Rn}}$}};
\node at (17.18,-5.18) {\tiny 86};
\node[fill=white, draw=black, circle, inner sep=0.4pt] at (17.8,-5.8) {\tiny \textbf{?}};
\fill[blue!18, rounded corners=0.7pt] (0,-6) -- (1,-6) -- (1,-7) -- cycle;
\fill[gray!8, rounded corners=0.7pt] (0,-6) -- (0,-7) -- (1,-7) -- cycle;
\draw[black, rounded corners=0.7pt] (0,-6) rectangle (1,-7);
\draw[black, dashed, very thin] (0,-6) -- (1,-7);
\node at (0.7,-6.25) {\tiny \textbf{Fr}};
\node at (0.28,-6.73) {\tiny \textbf{$\overline{\mathrm{Fr}}$}};
\node at (0.18,-6.18) {\tiny 87};
\node[fill=white, draw=black, circle, inner sep=0.4pt] at (0.8,-6.8) {\tiny \textbf{?}};
\fill[blue!18, rounded corners=0.7pt] (1,-6) -- (2,-6) -- (2,-7) -- cycle;
\fill[gray!8, rounded corners=0.7pt] (1,-6) -- (1,-7) -- (2,-7) -- cycle;
\draw[black, rounded corners=0.7pt] (1,-6) rectangle (2,-7);
\draw[black, dashed, very thin] (1,-6) -- (2,-7);
\node at (1.7,-6.25) {\tiny \textbf{Ra}};
\node at (1.28,-6.73) {\tiny \textbf{$\overline{\mathrm{Ra}}$}};
\node at (1.18,-6.18) {\tiny 88};
\node[fill=white, draw=black, circle, inner sep=0.4pt] at (1.8,-6.8) {\tiny \textbf{?}};
\fill[orange!20, rounded corners=0.7pt] (2,-6) -- (3,-6) -- (3,-7) -- cycle;
\fill[gray!8, rounded corners=0.7pt] (2,-6) -- (2,-7) -- (3,-7) -- cycle;
\draw[black, rounded corners=0.7pt] (2,-6) rectangle (3,-7);
\draw[black, dashed, very thin] (2,-6) -- (3,-7);
\node at (2.7,-6.25) {\tiny \textbf{Lr}};
\node at (2.2800000000000002,-6.73) {\tiny \textbf{$\overline{\mathrm{Lr}}$}};
\node at (2.18,-6.18) {\tiny 103};
\node[fill=white, draw=black, circle, inner sep=0.4pt] at (2.8,-6.8) {\tiny \textbf{?}};
\fill[orange!20, rounded corners=0.7pt] (3,-6) -- (4,-6) -- (4,-7) -- cycle;
\fill[gray!8, rounded corners=0.7pt] (3,-6) -- (3,-7) -- (4,-7) -- cycle;
\draw[black, rounded corners=0.7pt] (3,-6) rectangle (4,-7);
\draw[black, dashed, very thin] (3,-6) -- (4,-7);
\node at (3.7,-6.25) {\tiny \textbf{Rf}};
\node at (3.2800000000000002,-6.73) {\tiny \textbf{$\overline{\mathrm{Rf}}$}};
\node at (3.18,-6.18) {\tiny 104};
\node[fill=white, draw=black, circle, inner sep=0.4pt] at (3.8,-6.8) {\tiny \textbf{?}};
\fill[orange!20, rounded corners=0.7pt] (4,-6) -- (5,-6) -- (5,-7) -- cycle;
\fill[gray!8, rounded corners=0.7pt] (4,-6) -- (4,-7) -- (5,-7) -- cycle;
\draw[black, rounded corners=0.7pt] (4,-6) rectangle (5,-7);
\draw[black, dashed, very thin] (4,-6) -- (5,-7);
\node at (4.7,-6.25) {\tiny \textbf{Db}};
\node at (4.28,-6.73) {\tiny \textbf{$\overline{\mathrm{Db}}$}};
\node at (4.18,-6.18) {\tiny 105};
\node[fill=white, draw=black, circle, inner sep=0.4pt] at (4.8,-6.8) {\tiny \textbf{?}};
\fill[orange!20, rounded corners=0.7pt] (5,-6) -- (6,-6) -- (6,-7) -- cycle;
\fill[gray!8, rounded corners=0.7pt] (5,-6) -- (5,-7) -- (6,-7) -- cycle;
\draw[black, rounded corners=0.7pt] (5,-6) rectangle (6,-7);
\draw[black, dashed, very thin] (5,-6) -- (6,-7);
\node at (5.7,-6.25) {\tiny \textbf{Sg}};
\node at (5.28,-6.73) {\tiny \textbf{$\overline{\mathrm{Sg}}$}};
\node at (5.18,-6.18) {\tiny 106};
\node[fill=white, draw=black, circle, inner sep=0.4pt] at (5.8,-6.8) {\tiny \textbf{?}};
\fill[orange!20, rounded corners=0.7pt] (6,-6) -- (7,-6) -- (7,-7) -- cycle;
\fill[gray!8, rounded corners=0.7pt] (6,-6) -- (6,-7) -- (7,-7) -- cycle;
\draw[black, rounded corners=0.7pt] (6,-6) rectangle (7,-7);
\draw[black, dashed, very thin] (6,-6) -- (7,-7);
\node at (6.7,-6.25) {\tiny \textbf{Bh}};
\node at (6.28,-6.73) {\tiny \textbf{$\overline{\mathrm{Bh}}$}};
\node at (6.18,-6.18) {\tiny 107};
\node[fill=white, draw=black, circle, inner sep=0.4pt] at (6.8,-6.8) {\tiny \textbf{?}};
\fill[orange!20, rounded corners=0.7pt] (7,-6) -- (8,-6) -- (8,-7) -- cycle;
\fill[gray!8, rounded corners=0.7pt] (7,-6) -- (7,-7) -- (8,-7) -- cycle;
\draw[black, rounded corners=0.7pt] (7,-6) rectangle (8,-7);
\draw[black, dashed, very thin] (7,-6) -- (8,-7);
\node at (7.7,-6.25) {\tiny \textbf{Hs}};
\node at (7.28,-6.73) {\tiny \textbf{$\overline{\mathrm{Hs}}$}};
\node at (7.18,-6.18) {\tiny 108};
\node[fill=white, draw=black, circle, inner sep=0.4pt] at (7.8,-6.8) {\tiny \textbf{?}};
\fill[orange!20, rounded corners=0.7pt] (8,-6) -- (9,-6) -- (9,-7) -- cycle;
\fill[gray!8, rounded corners=0.7pt] (8,-6) -- (8,-7) -- (9,-7) -- cycle;
\draw[black, rounded corners=0.7pt] (8,-6) rectangle (9,-7);
\draw[black, dashed, very thin] (8,-6) -- (9,-7);
\node at (8.7,-6.25) {\tiny \textbf{Mt}};
\node at (8.28,-6.73) {\tiny \textbf{$\overline{\mathrm{Mt}}$}};
\node at (8.18,-6.18) {\tiny 109};
\node[fill=white, draw=black, circle, inner sep=0.4pt] at (8.8,-6.8) {\tiny \textbf{?}};
\fill[orange!20, rounded corners=0.7pt] (9,-6) -- (10,-6) -- (10,-7) -- cycle;
\fill[gray!8, rounded corners=0.7pt] (9,-6) -- (9,-7) -- (10,-7) -- cycle;
\draw[black, rounded corners=0.7pt] (9,-6) rectangle (10,-7);
\draw[black, dashed, very thin] (9,-6) -- (10,-7);
\node at (9.7,-6.25) {\tiny \textbf{Ds}};
\node at (9.28,-6.73) {\tiny \textbf{$\overline{\mathrm{Ds}}$}};
\node at (9.18,-6.18) {\tiny 110};
\node[fill=white, draw=black, circle, inner sep=0.4pt] at (9.8,-6.8) {\tiny \textbf{?}};
\fill[orange!20, rounded corners=0.7pt] (10,-6) -- (11,-6) -- (11,-7) -- cycle;
\fill[gray!8, rounded corners=0.7pt] (10,-6) -- (10,-7) -- (11,-7) -- cycle;
\draw[black, rounded corners=0.7pt] (10,-6) rectangle (11,-7);
\draw[black, dashed, very thin] (10,-6) -- (11,-7);
\node at (10.7,-6.25) {\tiny \textbf{Rg}};
\node at (10.28,-6.73) {\tiny \textbf{$\overline{\mathrm{Rg}}$}};
\node at (10.18,-6.18) {\tiny 111};
\node[fill=white, draw=black, circle, inner sep=0.4pt] at (10.8,-6.8) {\tiny \textbf{?}};
\fill[orange!20, rounded corners=0.7pt] (11,-6) -- (12,-6) -- (12,-7) -- cycle;
\fill[gray!8, rounded corners=0.7pt] (11,-6) -- (11,-7) -- (12,-7) -- cycle;
\draw[black, rounded corners=0.7pt] (11,-6) rectangle (12,-7);
\draw[black, dashed, very thin] (11,-6) -- (12,-7);
\node at (11.7,-6.25) {\tiny \textbf{Cn}};
\node at (11.28,-6.73) {\tiny \textbf{$\overline{\mathrm{Cn}}$}};
\node at (11.18,-6.18) {\tiny 112};
\node[fill=white, draw=black, circle, inner sep=0.4pt] at (11.8,-6.8) {\tiny \textbf{?}};
\fill[green!18, rounded corners=0.7pt] (12,-6) -- (13,-6) -- (13,-7) -- cycle;
\fill[gray!8, rounded corners=0.7pt] (12,-6) -- (12,-7) -- (13,-7) -- cycle;
\draw[black, rounded corners=0.7pt] (12,-6) rectangle (13,-7);
\draw[black, dashed, very thin] (12,-6) -- (13,-7);
\node at (12.7,-6.25) {\tiny \textbf{Nh}};
\node at (12.28,-6.73) {\tiny \textbf{$\overline{\mathrm{Nh}}$}};
\node at (12.18,-6.18) {\tiny 113};
\node[fill=white, draw=black, circle, inner sep=0.4pt] at (12.8,-6.8) {\tiny \textbf{?}};
\fill[green!18, rounded corners=0.7pt] (13,-6) -- (14,-6) -- (14,-7) -- cycle;
\fill[gray!8, rounded corners=0.7pt] (13,-6) -- (13,-7) -- (14,-7) -- cycle;
\draw[black, rounded corners=0.7pt] (13,-6) rectangle (14,-7);
\draw[black, dashed, very thin] (13,-6) -- (14,-7);
\node at (13.7,-6.25) {\tiny \textbf{Fl}};
\node at (13.28,-6.73) {\tiny \textbf{$\overline{\mathrm{Fl}}$}};
\node at (13.18,-6.18) {\tiny 114};
\node[fill=white, draw=black, circle, inner sep=0.4pt] at (13.8,-6.8) {\tiny \textbf{?}};
\fill[green!18, rounded corners=0.7pt] (14,-6) -- (15,-6) -- (15,-7) -- cycle;
\fill[gray!8, rounded corners=0.7pt] (14,-6) -- (14,-7) -- (15,-7) -- cycle;
\draw[black, rounded corners=0.7pt] (14,-6) rectangle (15,-7);
\draw[black, dashed, very thin] (14,-6) -- (15,-7);
\node at (14.7,-6.25) {\tiny \textbf{Mc}};
\node at (14.28,-6.73) {\tiny \textbf{$\overline{\mathrm{Mc}}$}};
\node at (14.18,-6.18) {\tiny 115};
\node[fill=white, draw=black, circle, inner sep=0.4pt] at (14.8,-6.8) {\tiny \textbf{?}};
\fill[green!18, rounded corners=0.7pt] (15,-6) -- (16,-6) -- (16,-7) -- cycle;
\fill[gray!8, rounded corners=0.7pt] (15,-6) -- (15,-7) -- (16,-7) -- cycle;
\draw[black, rounded corners=0.7pt] (15,-6) rectangle (16,-7);
\draw[black, dashed, very thin] (15,-6) -- (16,-7);
\node at (15.7,-6.25) {\tiny \textbf{Lv}};
\node at (15.28,-6.73) {\tiny \textbf{$\overline{\mathrm{Lv}}$}};
\node at (15.18,-6.18) {\tiny 116};
\node[fill=white, draw=black, circle, inner sep=0.4pt] at (15.8,-6.8) {\tiny \textbf{?}};
\fill[green!18, rounded corners=0.7pt] (16,-6) -- (17,-6) -- (17,-7) -- cycle;
\fill[gray!8, rounded corners=0.7pt] (16,-6) -- (16,-7) -- (17,-7) -- cycle;
\draw[black, rounded corners=0.7pt] (16,-6) rectangle (17,-7);
\draw[black, dashed, very thin] (16,-6) -- (17,-7);
\node at (16.7,-6.25) {\tiny \textbf{Ts}};
\node at (16.28,-6.73) {\tiny \textbf{$\overline{\mathrm{Ts}}$}};
\node at (16.18,-6.18) {\tiny 117};
\node[fill=white, draw=black, circle, inner sep=0.4pt] at (16.8,-6.8) {\tiny \textbf{?}};
\fill[green!18, rounded corners=0.7pt] (17,-6) -- (18,-6) -- (18,-7) -- cycle;
\fill[gray!8, rounded corners=0.7pt] (17,-6) -- (17,-7) -- (18,-7) -- cycle;
\draw[black, rounded corners=0.7pt] (17,-6) rectangle (18,-7);
\draw[black, dashed, very thin] (17,-6) -- (18,-7);
\node at (17.7,-6.25) {\tiny \textbf{Og}};
\node at (17.28,-6.73) {\tiny \textbf{$\overline{\mathrm{Og}}$}};
\node at (17.18,-6.18) {\tiny 118};
\node[fill=white, draw=black, circle, inner sep=0.4pt] at (17.8,-6.8) {\tiny \textbf{?}};
\fill[red!18, rounded corners=0.7pt] (3,-7.5) -- (4,-7.5) -- (4,-8.5) -- cycle;
\fill[gray!8, rounded corners=0.7pt] (3,-7.5) -- (3,-8.5) -- (4,-8.5) -- cycle;
\draw[black, rounded corners=0.7pt] (3,-7.5) rectangle (4,-8.5);
\draw[black, dashed, very thin] (3,-7.5) -- (4,-8.5);
\node at (3.7,-7.75) {\tiny \textbf{La}};
\node at (3.2800000000000002,-8.23) {\tiny \textbf{$\overline{\mathrm{La}}$}};
\node at (3.18,-7.68) {\tiny 57};
\node[fill=white, draw=black, circle, inner sep=0.4pt] at (3.8,-8.3) {\tiny \textbf{?}};
\fill[red!18, rounded corners=0.7pt] (4,-7.5) -- (5,-7.5) -- (5,-8.5) -- cycle;
\fill[gray!8, rounded corners=0.7pt] (4,-7.5) -- (4,-8.5) -- (5,-8.5) -- cycle;
\draw[black, rounded corners=0.7pt] (4,-7.5) rectangle (5,-8.5);
\draw[black, dashed, very thin] (4,-7.5) -- (5,-8.5);
\node at (4.7,-7.75) {\tiny \textbf{Ce}};
\node at (4.28,-8.23) {\tiny \textbf{$\overline{\mathrm{Ce}}$}};
\node at (4.18,-7.68) {\tiny 58};
\node[fill=white, draw=black, circle, inner sep=0.4pt] at (4.8,-8.3) {\tiny \textbf{?}};
\fill[red!18, rounded corners=0.7pt] (5,-7.5) -- (6,-7.5) -- (6,-8.5) -- cycle;
\fill[gray!8, rounded corners=0.7pt] (5,-7.5) -- (5,-8.5) -- (6,-8.5) -- cycle;
\draw[black, rounded corners=0.7pt] (5,-7.5) rectangle (6,-8.5);
\draw[black, dashed, very thin] (5,-7.5) -- (6,-8.5);
\node at (5.7,-7.75) {\tiny \textbf{Pr}};
\node at (5.28,-8.23) {\tiny \textbf{$\overline{\mathrm{Pr}}$}};
\node at (5.18,-7.68) {\tiny 59};
\node[fill=white, draw=black, circle, inner sep=0.4pt] at (5.8,-8.3) {\tiny \textbf{?}};
\fill[red!18, rounded corners=0.7pt] (6,-7.5) -- (7,-7.5) -- (7,-8.5) -- cycle;
\fill[gray!8, rounded corners=0.7pt] (6,-7.5) -- (6,-8.5) -- (7,-8.5) -- cycle;
\draw[black, rounded corners=0.7pt] (6,-7.5) rectangle (7,-8.5);
\draw[black, dashed, very thin] (6,-7.5) -- (7,-8.5);
\node at (6.7,-7.75) {\tiny \textbf{Nd}};
\node at (6.28,-8.23) {\tiny \textbf{$\overline{\mathrm{Nd}}$}};
\node at (6.18,-7.68) {\tiny 60};
\node[fill=white, draw=black, circle, inner sep=0.4pt] at (6.8,-8.3) {\tiny \textbf{?}};
\fill[red!18, rounded corners=0.7pt] (7,-7.5) -- (8,-7.5) -- (8,-8.5) -- cycle;
\fill[gray!8, rounded corners=0.7pt] (7,-7.5) -- (7,-8.5) -- (8,-8.5) -- cycle;
\draw[black, rounded corners=0.7pt] (7,-7.5) rectangle (8,-8.5);
\draw[black, dashed, very thin] (7,-7.5) -- (8,-8.5);
\node at (7.7,-7.75) {\tiny \textbf{Pm}};
\node at (7.28,-8.23) {\tiny \textbf{$\overline{\mathrm{Pm}}$}};
\node at (7.18,-7.68) {\tiny 61};
\node[fill=white, draw=black, circle, inner sep=0.4pt] at (7.8,-8.3) {\tiny \textbf{?}};
\fill[red!18, rounded corners=0.7pt] (8,-7.5) -- (9,-7.5) -- (9,-8.5) -- cycle;
\fill[gray!8, rounded corners=0.7pt] (8,-7.5) -- (8,-8.5) -- (9,-8.5) -- cycle;
\draw[black, rounded corners=0.7pt] (8,-7.5) rectangle (9,-8.5);
\draw[black, dashed, very thin] (8,-7.5) -- (9,-8.5);
\node at (8.7,-7.75) {\tiny \textbf{Sm}};
\node at (8.28,-8.23) {\tiny \textbf{$\overline{\mathrm{Sm}}$}};
\node at (8.18,-7.68) {\tiny 62};
\node[fill=white, draw=black, circle, inner sep=0.4pt] at (8.8,-8.3) {\tiny \textbf{?}};
\fill[red!18, rounded corners=0.7pt] (9,-7.5) -- (10,-7.5) -- (10,-8.5) -- cycle;
\fill[gray!8, rounded corners=0.7pt] (9,-7.5) -- (9,-8.5) -- (10,-8.5) -- cycle;
\draw[black, rounded corners=0.7pt] (9,-7.5) rectangle (10,-8.5);
\draw[black, dashed, very thin] (9,-7.5) -- (10,-8.5);
\node at (9.7,-7.75) {\tiny \textbf{Eu}};
\node at (9.28,-8.23) {\tiny \textbf{$\overline{\mathrm{Eu}}$}};
\node at (9.18,-7.68) {\tiny 63};
\node[fill=white, draw=black, circle, inner sep=0.4pt] at (9.8,-8.3) {\tiny \textbf{?}};
\fill[red!18, rounded corners=0.7pt] (10,-7.5) -- (11,-7.5) -- (11,-8.5) -- cycle;
\fill[gray!8, rounded corners=0.7pt] (10,-7.5) -- (10,-8.5) -- (11,-8.5) -- cycle;
\draw[black, rounded corners=0.7pt] (10,-7.5) rectangle (11,-8.5);
\draw[black, dashed, very thin] (10,-7.5) -- (11,-8.5);
\node at (10.7,-7.75) {\tiny \textbf{Gd}};
\node at (10.28,-8.23) {\tiny \textbf{$\overline{\mathrm{Gd}}$}};
\node at (10.18,-7.68) {\tiny 64};
\node[fill=white, draw=black, circle, inner sep=0.4pt] at (10.8,-8.3) {\tiny \textbf{?}};
\fill[red!18, rounded corners=0.7pt] (11,-7.5) -- (12,-7.5) -- (12,-8.5) -- cycle;
\fill[gray!8, rounded corners=0.7pt] (11,-7.5) -- (11,-8.5) -- (12,-8.5) -- cycle;
\draw[black, rounded corners=0.7pt] (11,-7.5) rectangle (12,-8.5);
\draw[black, dashed, very thin] (11,-7.5) -- (12,-8.5);
\node at (11.7,-7.75) {\tiny \textbf{Tb}};
\node at (11.28,-8.23) {\tiny \textbf{$\overline{\mathrm{Tb}}$}};
\node at (11.18,-7.68) {\tiny 65};
\node[fill=white, draw=black, circle, inner sep=0.4pt] at (11.8,-8.3) {\tiny \textbf{?}};
\fill[green!18, rounded corners=0.7pt] (12,-7.5) -- (13,-7.5) -- (13,-8.5) -- cycle;
\fill[gray!8, rounded corners=0.7pt] (12,-7.5) -- (12,-8.5) -- (13,-8.5) -- cycle;
\draw[black, rounded corners=0.7pt] (12,-7.5) rectangle (13,-8.5);
\draw[black, dashed, very thin] (12,-7.5) -- (13,-8.5);
\node at (12.7,-7.75) {\tiny \textbf{Dy}};
\node at (12.28,-8.23) {\tiny \textbf{$\overline{\mathrm{Dy}}$}};
\node at (12.18,-7.68) {\tiny 66};
\node[fill=white, draw=black, circle, inner sep=0.4pt] at (12.8,-8.3) {\tiny \textbf{?}};
\fill[green!18, rounded corners=0.7pt] (13,-7.5) -- (14,-7.5) -- (14,-8.5) -- cycle;
\fill[gray!8, rounded corners=0.7pt] (13,-7.5) -- (13,-8.5) -- (14,-8.5) -- cycle;
\draw[black, rounded corners=0.7pt] (13,-7.5) rectangle (14,-8.5);
\draw[black, dashed, very thin] (13,-7.5) -- (14,-8.5);
\node at (13.7,-7.75) {\tiny \textbf{Ho}};
\node at (13.28,-8.23) {\tiny \textbf{$\overline{\mathrm{Ho}}$}};
\node at (13.18,-7.68) {\tiny 67};
\node[fill=white, draw=black, circle, inner sep=0.4pt] at (13.8,-8.3) {\tiny \textbf{?}};
\fill[green!18, rounded corners=0.7pt] (14,-7.5) -- (15,-7.5) -- (15,-8.5) -- cycle;
\fill[gray!8, rounded corners=0.7pt] (14,-7.5) -- (14,-8.5) -- (15,-8.5) -- cycle;
\draw[black, rounded corners=0.7pt] (14,-7.5) rectangle (15,-8.5);
\draw[black, dashed, very thin] (14,-7.5) -- (15,-8.5);
\node at (14.7,-7.75) {\tiny \textbf{Er}};
\node at (14.28,-8.23) {\tiny \textbf{$\overline{\mathrm{Er}}$}};
\node at (14.18,-7.68) {\tiny 68};
\node[fill=white, draw=black, circle, inner sep=0.4pt] at (14.8,-8.3) {\tiny \textbf{?}};
\fill[green!18, rounded corners=0.7pt] (15,-7.5) -- (16,-7.5) -- (16,-8.5) -- cycle;
\fill[gray!8, rounded corners=0.7pt] (15,-7.5) -- (15,-8.5) -- (16,-8.5) -- cycle;
\draw[black, rounded corners=0.7pt] (15,-7.5) rectangle (16,-8.5);
\draw[black, dashed, very thin] (15,-7.5) -- (16,-8.5);
\node at (15.7,-7.75) {\tiny \textbf{Tm}};
\node at (15.28,-8.23) {\tiny \textbf{$\overline{\mathrm{Tm}}$}};
\node at (15.18,-7.68) {\tiny 69};
\node[fill=white, draw=black, circle, inner sep=0.4pt] at (15.8,-8.3) {\tiny \textbf{?}};
\fill[green!18, rounded corners=0.7pt] (16,-7.5) -- (17,-7.5) -- (17,-8.5) -- cycle;
\fill[gray!8, rounded corners=0.7pt] (16,-7.5) -- (16,-8.5) -- (17,-8.5) -- cycle;
\draw[black, rounded corners=0.7pt] (16,-7.5) rectangle (17,-8.5);
\draw[black, dashed, very thin] (16,-7.5) -- (17,-8.5);
\node at (16.7,-7.75) {\tiny \textbf{Yb}};
\node at (16.28,-8.23) {\tiny \textbf{$\overline{\mathrm{Yb}}$}};
\node at (16.18,-7.68) {\tiny 70};
\node[fill=white, draw=black, circle, inner sep=0.4pt] at (16.8,-8.3) {\tiny \textbf{?}};
\fill[red!18, rounded corners=0.7pt] (3,-8.5) -- (4,-8.5) -- (4,-9.5) -- cycle;
\fill[gray!8, rounded corners=0.7pt] (3,-8.5) -- (3,-9.5) -- (4,-9.5) -- cycle;
\draw[black, rounded corners=0.7pt] (3,-8.5) rectangle (4,-9.5);
\draw[black, dashed, very thin] (3,-8.5) -- (4,-9.5);
\node at (3.7,-8.75) {\tiny \textbf{Ac}};
\node at (3.2800000000000002,-9.23) {\tiny \textbf{$\overline{\mathrm{Ac}}$}};
\node at (3.18,-8.68) {\tiny 89};
\node[fill=white, draw=black, circle, inner sep=0.4pt] at (3.8,-9.3) {\tiny \textbf{?}};
\fill[red!18, rounded corners=0.7pt] (4,-8.5) -- (5,-8.5) -- (5,-9.5) -- cycle;
\fill[gray!8, rounded corners=0.7pt] (4,-8.5) -- (4,-9.5) -- (5,-9.5) -- cycle;
\draw[black, rounded corners=0.7pt] (4,-8.5) rectangle (5,-9.5);
\draw[black, dashed, very thin] (4,-8.5) -- (5,-9.5);
\node at (4.7,-8.75) {\tiny \textbf{Th}};
\node at (4.28,-9.23) {\tiny \textbf{$\overline{\mathrm{Th}}$}};
\node at (4.18,-8.68) {\tiny 90};
\node[fill=white, draw=black, circle, inner sep=0.4pt] at (4.8,-9.3) {\tiny \textbf{?}};
\fill[red!18, rounded corners=0.7pt] (5,-8.5) -- (6,-8.5) -- (6,-9.5) -- cycle;
\fill[gray!8, rounded corners=0.7pt] (5,-8.5) -- (5,-9.5) -- (6,-9.5) -- cycle;
\draw[black, rounded corners=0.7pt] (5,-8.5) rectangle (6,-9.5);
\draw[black, dashed, very thin] (5,-8.5) -- (6,-9.5);
\node at (5.7,-8.75) {\tiny \textbf{Pa}};
\node at (5.28,-9.23) {\tiny \textbf{$\overline{\mathrm{Pa}}$}};
\node at (5.18,-8.68) {\tiny 91};
\node[fill=white, draw=black, circle, inner sep=0.4pt] at (5.8,-9.3) {\tiny \textbf{?}};
\fill[red!18, rounded corners=0.7pt] (6,-8.5) -- (7,-8.5) -- (7,-9.5) -- cycle;
\fill[gray!8, rounded corners=0.7pt] (6,-8.5) -- (6,-9.5) -- (7,-9.5) -- cycle;
\draw[black, rounded corners=0.7pt] (6,-8.5) rectangle (7,-9.5);
\draw[black, dashed, very thin] (6,-8.5) -- (7,-9.5);
\node at (6.7,-8.75) {\tiny \textbf{U}};
\node at (6.28,-9.23) {\tiny \textbf{$\overline{\mathrm{U}}$}};
\node at (6.18,-8.68) {\tiny 92};
\node[fill=white, draw=black, circle, inner sep=0.4pt] at (6.8,-9.3) {\tiny \textbf{?}};
\fill[red!18, rounded corners=0.7pt] (7,-8.5) -- (8,-8.5) -- (8,-9.5) -- cycle;
\fill[gray!8, rounded corners=0.7pt] (7,-8.5) -- (7,-9.5) -- (8,-9.5) -- cycle;
\draw[black, rounded corners=0.7pt] (7,-8.5) rectangle (8,-9.5);
\draw[black, dashed, very thin] (7,-8.5) -- (8,-9.5);
\node at (7.7,-8.75) {\tiny \textbf{Np}};
\node at (7.28,-9.23) {\tiny \textbf{$\overline{\mathrm{Np}}$}};
\node at (7.18,-8.68) {\tiny 93};
\node[fill=white, draw=black, circle, inner sep=0.4pt] at (7.8,-9.3) {\tiny \textbf{?}};
\fill[red!18, rounded corners=0.7pt] (8,-8.5) -- (9,-8.5) -- (9,-9.5) -- cycle;
\fill[gray!8, rounded corners=0.7pt] (8,-8.5) -- (8,-9.5) -- (9,-9.5) -- cycle;
\draw[black, rounded corners=0.7pt] (8,-8.5) rectangle (9,-9.5);
\draw[black, dashed, very thin] (8,-8.5) -- (9,-9.5);
\node at (8.7,-8.75) {\tiny \textbf{Pu}};
\node at (8.28,-9.23) {\tiny \textbf{$\overline{\mathrm{Pu}}$}};
\node at (8.18,-8.68) {\tiny 94};
\node[fill=white, draw=black, circle, inner sep=0.4pt] at (8.8,-9.3) {\tiny \textbf{?}};
\fill[red!18, rounded corners=0.7pt] (9,-8.5) -- (10,-8.5) -- (10,-9.5) -- cycle;
\fill[gray!8, rounded corners=0.7pt] (9,-8.5) -- (9,-9.5) -- (10,-9.5) -- cycle;
\draw[black, rounded corners=0.7pt] (9,-8.5) rectangle (10,-9.5);
\draw[black, dashed, very thin] (9,-8.5) -- (10,-9.5);
\node at (9.7,-8.75) {\tiny \textbf{Am}};
\node at (9.28,-9.23) {\tiny \textbf{$\overline{\mathrm{Am}}$}};
\node at (9.18,-8.68) {\tiny 95};
\node[fill=white, draw=black, circle, inner sep=0.4pt] at (9.8,-9.3) {\tiny \textbf{?}};
\fill[red!18, rounded corners=0.7pt] (10,-8.5) -- (11,-8.5) -- (11,-9.5) -- cycle;
\fill[gray!8, rounded corners=0.7pt] (10,-8.5) -- (10,-9.5) -- (11,-9.5) -- cycle;
\draw[black, rounded corners=0.7pt] (10,-8.5) rectangle (11,-9.5);
\draw[black, dashed, very thin] (10,-8.5) -- (11,-9.5);
\node at (10.7,-8.75) {\tiny \textbf{Cm}};
\node at (10.28,-9.23) {\tiny \textbf{$\overline{\mathrm{Cm}}$}};
\node at (10.18,-8.68) {\tiny 96};
\node[fill=white, draw=black, circle, inner sep=0.4pt] at (10.8,-9.3) {\tiny \textbf{?}};
\fill[red!18, rounded corners=0.7pt] (11,-8.5) -- (12,-8.5) -- (12,-9.5) -- cycle;
\fill[gray!8, rounded corners=0.7pt] (11,-8.5) -- (11,-9.5) -- (12,-9.5) -- cycle;
\draw[black, rounded corners=0.7pt] (11,-8.5) rectangle (12,-9.5);
\draw[black, dashed, very thin] (11,-8.5) -- (12,-9.5);
\node at (11.7,-8.75) {\tiny \textbf{Bk}};
\node at (11.28,-9.23) {\tiny \textbf{$\overline{\mathrm{Bk}}$}};
\node at (11.18,-8.68) {\tiny 97};
\node[fill=white, draw=black, circle, inner sep=0.4pt] at (11.8,-9.3) {\tiny \textbf{?}};
\fill[green!18, rounded corners=0.7pt] (12,-8.5) -- (13,-8.5) -- (13,-9.5) -- cycle;
\fill[gray!8, rounded corners=0.7pt] (12,-8.5) -- (12,-9.5) -- (13,-9.5) -- cycle;
\draw[black, rounded corners=0.7pt] (12,-8.5) rectangle (13,-9.5);
\draw[black, dashed, very thin] (12,-8.5) -- (13,-9.5);
\node at (12.7,-8.75) {\tiny \textbf{Cf}};
\node at (12.28,-9.23) {\tiny \textbf{$\overline{\mathrm{Cf}}$}};
\node at (12.18,-8.68) {\tiny 98};
\node[fill=white, draw=black, circle, inner sep=0.4pt] at (12.8,-9.3) {\tiny \textbf{?}};
\fill[green!18, rounded corners=0.7pt] (13,-8.5) -- (14,-8.5) -- (14,-9.5) -- cycle;
\fill[gray!8, rounded corners=0.7pt] (13,-8.5) -- (13,-9.5) -- (14,-9.5) -- cycle;
\draw[black, rounded corners=0.7pt] (13,-8.5) rectangle (14,-9.5);
\draw[black, dashed, very thin] (13,-8.5) -- (14,-9.5);
\node at (13.7,-8.75) {\tiny \textbf{Es}};
\node at (13.28,-9.23) {\tiny \textbf{$\overline{\mathrm{Es}}$}};
\node at (13.18,-8.68) {\tiny 99};
\node[fill=white, draw=black, circle, inner sep=0.4pt] at (13.8,-9.3) {\tiny \textbf{?}};
\fill[green!18, rounded corners=0.7pt] (14,-8.5) -- (15,-8.5) -- (15,-9.5) -- cycle;
\fill[gray!8, rounded corners=0.7pt] (14,-8.5) -- (14,-9.5) -- (15,-9.5) -- cycle;
\draw[black, rounded corners=0.7pt] (14,-8.5) rectangle (15,-9.5);
\draw[black, dashed, very thin] (14,-8.5) -- (15,-9.5);
\node at (14.7,-8.75) {\tiny \textbf{Fm}};
\node at (14.28,-9.23) {\tiny \textbf{$\overline{\mathrm{Fm}}$}};
\node at (14.18,-8.68) {\tiny 100};
\node[fill=white, draw=black, circle, inner sep=0.4pt] at (14.8,-9.3) {\tiny \textbf{?}};
\fill[green!18, rounded corners=0.7pt] (15,-8.5) -- (16,-8.5) -- (16,-9.5) -- cycle;
\fill[gray!8, rounded corners=0.7pt] (15,-8.5) -- (15,-9.5) -- (16,-9.5) -- cycle;
\draw[black, rounded corners=0.7pt] (15,-8.5) rectangle (16,-9.5);
\draw[black, dashed, very thin] (15,-8.5) -- (16,-9.5);
\node at (15.7,-8.75) {\tiny \textbf{Md}};
\node at (15.28,-9.23) {\tiny \textbf{$\overline{\mathrm{Md}}$}};
\node at (15.18,-8.68) {\tiny 101};
\node[fill=white, draw=black, circle, inner sep=0.4pt] at (15.8,-9.3) {\tiny \textbf{?}};
\fill[green!18, rounded corners=0.7pt] (16,-8.5) -- (17,-8.5) -- (17,-9.5) -- cycle;
\fill[gray!8, rounded corners=0.7pt] (16,-8.5) -- (16,-9.5) -- (17,-9.5) -- cycle;
\draw[black, rounded corners=0.7pt] (16,-8.5) rectangle (17,-9.5);
\draw[black, dashed, very thin] (16,-8.5) -- (17,-9.5);
\node at (16.7,-8.75) {\tiny \textbf{No}};
\node at (16.28,-9.23) {\tiny \textbf{$\overline{\mathrm{No}}$}};
\node at (16.18,-8.68) {\tiny 102};
\node[fill=white, draw=black, circle, inner sep=0.4pt] at (16.8,-9.3) {\tiny \textbf{?}};

\node[anchor=west] at (0.0,-10.309999999999999) {\scriptsize \textbf{Block:}};
\draw[fill=blue!18, draw=black] (1.65,-10.2) rectangle (2.2,-10.5);
\node[anchor=west] at (2.3499999999999996,-10.35) {\scriptsize s};
\draw[fill=green!18, draw=black] (2.8499999999999996,-10.2) rectangle (3.3999999999999995,-10.5);
\node[anchor=west] at (3.55,-10.35) {\scriptsize p};
\draw[fill=orange!20, draw=black] (4.05,-10.2) rectangle (4.6,-10.5);
\node[anchor=west] at (4.75,-10.35) {\scriptsize d};
\draw[fill=red!18, draw=black] (5.25,-10.2) rectangle (5.8,-10.5);
\node[anchor=west] at (5.95,-10.35) {\scriptsize f};
\node[anchor=west] at (7.4,-10.35) {\scriptsize Cell size: $0.68\,\mathrm{cm}\times0.68\,\mathrm{cm}$.};
\node[anchor=west] at (0.0,-11.009999999999998) {\scriptsize \textbf{Marker:}};
\node[anchor=west] at (2.2,-11.049999999999999) {\scriptsize $=$ CPT-compatible;\quad K antinucleus;\quad ? open;\quad R residual.};
\node[anchor=west] at (0.0,-11.709999999999999) {\scriptsize \textbf{Anti status:}};
\draw[fill=green!60, draw=black] (2.9,-11.6) rectangle (3.45,-11.9);
\node[anchor=west] at (3.6,-11.75) {\scriptsize Atom};
\draw[fill=orange!60, draw=black] (5.1,-11.6) rectangle (5.65,-11.9);
\node[anchor=west] at (5.8,-11.75) {\scriptsize Nucleus};
\draw[fill=gray!8, draw=black] (7.6,-11.6) rectangle (8.15,-11.9);
\node[anchor=west] at (8.3,-11.75) {\scriptsize Open};
\end{tikzpicture}
\caption{Split-cell antimatter periodic-table framework. Classical positions follow the IUPAC source basis; antimatter markers are evidence classifications from the cited status table: measured CPT compatibility ($=$), antinucleus-only status (K), open neutral anti-atom status (?), or residual analysis (R).}
\label{fig:ptable-split}
\end{figure}


\section{Neuheitswert des Vorschlags}

Der Vorschlag ist nur dann stark, wenn er praezise abgegrenzt wird.

\subsection{Nicht neu}

Nicht neu sind folgende Punkte:

\begin{itemize}
  \item Antiatome als ladungskonjugierte Gegenstuecke normaler Atome.
  \item Die Erwartung identischer Spektren unter CPT.
  \item Die Notation mit Ueberstrich fuer Antiteilchen bzw. Antisysteme.
  \item Der Begriff Antiwasserstoff.
\end{itemize}

\subsection{Potenziell neuer Beitrag}

Neu und verteidigungsfaehig ist dagegen die Kombination folgender Elemente:

\begin{enumerate}
  \item \textbf{Explizite Trennung von Struktur und Evidenz:} Das Periodengesetz bleibt strukturell erhalten; die Tabelle visualisiert, welche Anti-Systeme experimentell realisiert sind.
  \item \textbf{Schichtenmodell statt Einzelfarbe:} Antikernstatus, neutraler Anti-Atomstatus, Spektroskopie, Chemie und Astrophysik werden getrennt.
  \item \textbf{Nomenklaturdisziplin:} Keine Erfindung neuer IUPAC-Namen, sondern Arbeitsnotation mit Anti-Praefix und Ueberstrich.
  \item \textbf{Preparedness-Charakter:} Die Tabelle ist eine Roadmap fuer zukuenftige Anti-Ionen, Anti-Molekuele und moegliche kosmische Antimaterie-Signaturen.
\end{enumerate}

Eine geeignete Claim-Formulierung fuer das Paper waere daher:

\begin{quote}
Wir schlagen kein neues Periodengesetz fuer Antimaterie vor, sondern ein statuscodiertes, CPT-isomorphes Preparedness Framework, das die klassische Orbitalstruktur mit dem experimentellen Realisierungsstand von Anti-Atomen, Antikernen, Anti-Ionen, Anti-Molekuelen und astrophysikalischen Antimaterie-Hypothesen verbindet.
\end{quote}

\section{Schlussfolgerung}

Ein Antimaterie-Periodensystem ist als eigenstaendiges chemisches Periodengesetz redundant, aber als epistemische Forschungslandkarte sinnvoll. Die wissenschaftlich robuste Position lautet:

\begin{itemize}
  \item Die Struktur des Antiperiodensystems ist nach CPT und QED isomorph zum klassischen Periodensystem.
  \item Die Namen der Antielemente sind nicht als eigene IUPAC-Reihe vergeben; empfohlen wird Anti-Praefix plus Ueberstrichsymbol.
  \item Antiwasserstoff ist der einzige neutrale Anti-Atomfall mit hoher experimenteller Charakterisierung.
  \item Antihelium ist bisher als Antikern, nicht als neutrales Antihelium-Atom gesichert.
  \item Anti-Molekuele, komplexe Antimaterie-Chemie und kosmische Antimaterie-Domaenen sind offene Forschungsfelder.
\end{itemize}

Damit ist der Vorschlag wissenschaftlich am staerksten, wenn er als \emph{Preparedness Framework} eingereicht wird: nicht als Behauptung einer anderen Antimaterie-Chemie, sondern als methodische, visuelle und nomenklatorische Vorbereitung auf zukuenftige experimentelle Systeme.

\begin{thebibliography}{99}

\bibitem{Dirac1928}
P. A. M. Dirac, \enquote{The quantum theory of the electron}, \emph{Proceedings of the Royal Society A} \textbf{117}, 610--624 (1928). DOI: 10.1098/rspa.1928.0023.

\bibitem{IUPACatomicnumber}
IUPAC Gold Book, \enquote{atomic number (proton number), Z}. DOI: 10.1351/goldbook.A00499.

\bibitem{IUPACpositron}
IUPAC Gold Book, \enquote{positron}. DOI: 10.1351/goldbook.P04769.

\bibitem{IUPACannihilation}
IUPAC Gold Book, \enquote{annihilation}. DOI: 10.1351/goldbook.A00366.

\bibitem{IUPAC113118}
IUPAC, \enquote{IUPAC Announces the Names of the Elements 113, 115, 117, and 118} (2016).

\bibitem{CERNstoring}
CERN, \enquote{Storing antihydrogen}. CERN Science/Physics/Antimatter, accessed 2026.

\bibitem{CERNantimatter}
CERN, \enquote{Antimatter}. CERN Science/Physics/Antimatter, accessed 2026.

\bibitem{CERNGBAR}
CERN, \enquote{GBAR: Gravitational Behaviour of Antihydrogen at Rest}. CERN Science/Experiments, accessed 2026.

\bibitem{ALPHA2018}
M. Ahmadi et al. (ALPHA Collaboration), \enquote{Characterization of the 1S--2S transition in antihydrogen}, \emph{Nature} \textbf{557}, 71--75 (2018). DOI: 10.1038/s41586-018-0017-2.

\bibitem{Anderson2023}
E. K. Anderson et al. (ALPHA Collaboration), \enquote{Observation of the effect of gravity on the motion of antimatter}, \emph{Nature} \textbf{621}, 716--722 (2023). DOI: 10.1038/s41586-023-06527-1.

\bibitem{Akbari2025}
R. Akbari et al. (ALPHA Collaboration), \enquote{Be$^{+}$ assisted, simultaneous confinement of more than 15000 antihydrogen atoms}, \emph{Nature Communications} \textbf{16}, 10106 (2025). DOI: 10.1038/s41467-025-65085-4.

\bibitem{Zammit2025}
M. C. Zammit, C. J. Baker, S. Jonsell, S. Eriksson and M. Charlton, \enquote{Antihydrogen chemistry}, \emph{Physical Review A} \textbf{111}, 050101 (2025). DOI: 10.1103/PhysRevA.111.050101.

\bibitem{STAR2011}
STAR Collaboration, \enquote{Observation of the antimatter helium-4 nucleus}, \emph{Nature} \textbf{473}, 353--356 (2011). DOI: 10.1038/nature10079.

\bibitem{STAR2024}
STAR Collaboration, \enquote{Observation of the antimatter hypernucleus $\,{}_{\bar\Lambda}^{4}\overline{\mathrm{H}}$}, \emph{Nature} \textbf{632}, 1026--1031 (2024). DOI: 10.1038/s41586-024-07823-0.

\bibitem{Hori2011}
M. Hori et al., \enquote{Two-photon laser spectroscopy of antiprotonic helium and the antiproton-to-electron mass ratio}, \emph{Nature} \textbf{475}, 484--488 (2011). DOI: 10.1038/nature10260.

\end{thebibliography}

\end{document}
